% !TEX program = lualatex
% ---------------------------------------------------------------------------------------------
%  Readback note 38: From amplitude sums to local correlations: filling levels at finitely many places, and the filling data of the admissible data
%
%  Build:  python3 tools/rb.py note <slug>
%
%  One unit of a Lean artifact, rewritten as ordinary mathematics. The prose is hand-written. Lean names,
%  line ranges and counts come from data.tex; quoted Lean statements are copied from the artifact by the build.
%
%  Rules of the text: standard terminology; every symbol is explained where it first appears and listed again
%  in the notation at the end; no letter has two meanings; a proof is complete or says what it takes for granted;
%  Lean names appear only in the LEAN lines and in the concordance.
%
%  The signature rule (the build enforces it): a function, relation or predicate that the artifact defines is
%  written only through its notation macro, which the build makes from the spec (`notation`) with one slot for
%  each explicit argument of the Lean definition. No argument is fixed and dropped, also not "for brevity" and
%  not for one section. Every argument and every hypothesis of a stated declaration stands inside its
%  statement, or in a Setting that the statement cites by number; only implicit and instance arguments (the
%  ambient objects) may stand in the conventions.
% ---------------------------------------------------------------------------------------------
\documentclass[11pt,a4paper]{article}
\ifdefined\pdfvariable\pdfvariable suppressoptionalinfo \numexpr 1+2+4+8\relax\fi
\usepackage[a4paper,hmargin=35mm,top=22mm,bottom=27mm,footskip=12mm,nohead]{geometry}

\usepackage{fontspec}
\directlua{
  luaotfload.add_fallback("gfallback", {"STIXTwoMath-Regular.otf:mode=node;", "DejaVuSans.ttf:mode=node;"})
  luaotfload.add_fallback("gmonofallback", {"DejaVuSansMono.ttf:mode=node;", "STIXTwoMath-Regular.otf:mode=node;"})
}
\setmainfont{NewCM10}[Extension=.otf,UprightFont=*-Book,ItalicFont=*-BookItalic,
  BoldFont=*-Bold,BoldItalicFont=*-BoldItalic,SmallCapsFont=*-Book,SmallCapsFeatures={Letters=SmallCaps},
  RawFeature={fallback=gfallback}]
\setsansfont{NewCMSans10}[Extension=.otf,UprightFont=*-Book,ItalicFont=*-BookOblique,
  BoldFont=*-Bold,BoldItalicFont=*-BoldOblique,RawFeature={fallback=gfallback}]
\setmonofont{JetBrainsMono}[Extension=.otf,UprightFont=*-Regular,ItalicFont=*-Italic,
  BoldFont=*-SemiBold,BoldItalicFont=*-SemiBoldItalic,Scale=MatchLowercase,
  RawFeature={-calt;-liga;fallback=gmonofallback}]
\usepackage{amsmath}
\usepackage{amsthm}
\usepackage[math-style=TeX]{unicode-math}
\setmathfont{NewCMMath-Book.otf}
\emergencystretch=2.5em\relax

\usepackage{xcolor}
\usepackage{tikz}
\usetikzlibrary{calc,arrows.meta,positioning,backgrounds}
\usepackage[protrusion=true,expansion=true]{microtype}
\usepackage{enumitem}
\usepackage{needspace}
\usepackage{booktabs}
\usepackage{array}
\usepackage{longtable}
\usepackage{fancyhdr}
\usepackage{lastpage}
\usepackage{multicol}
\usepackage[hidelinks,pdfusetitle]{hyperref}

\definecolor{ink}{HTML}{0F172A}
\definecolor{inksoft}{HTML}{475569}
\definecolor{inkfaint}{HTML}{94A3B8}
\definecolor{rulegrey}{HTML}{CBD5E1}
\definecolor{arith}{HTML}{C2410C}
\definecolor{est}{HTML}{7C3AED}
\definecolor{ana}{HTML}{0891B2}
\definecolor{ok}{HTML}{047857}
\definecolor{leancol}{HTML}{3730A3}

\title{From amplitude sums to local correlations: filling levels at finitely many places, and the filling data of the admissible data}
\author{grh-compact}
\date{}

% ------------------------------------------------------------------------------------- macros
\newcommand{\leanfont}{\ttfamily\fontsize{8.6}{10.5}\selectfont\color{leancol}}
\newcommand{\lean}[1]{\csname LN@#1\endcsname}
\newcommand{\fillin}[1]{{\color{arith}\itshape[#1]}}  % a place of the template that is not written yet; the build reports every one that is left
\newcommand{\leanquote}[1]{\csname LQ@#1\endcsname}
% generated by rb.py note; do not edit
\newcommand{\ConcDecls}{34}
\newcommand{\Pin}{72fead10}
\newcommand{\PinText}{that of commit {\leanfont 72fead10}}
\newcommand{\ArtifactName}{grh-compact}
\expandafter\newcommand\csname LN@relIndexlatStabilizer\endcsname{\href{https://github.com/lsernau/grh-lean/blob/72fead102f663d5c7a407964a13dc0cf9d977207/grh-compact/src/New/Filling/FillingAdmissible.lean\#L16-L28}{{\leanfont Data.\allowbreak{}relIndex\_latStabilizer}}}
\expandafter\newcommand\csname LN@cardclassesfilling\endcsname{\href{https://github.com/lsernau/grh-lean/blob/72fead102f663d5c7a407964a13dc0cf9d977207/grh-compact/src/New/Filling/FillingAdmissible.lean\#L30-L37}{{\leanfont Data.\allowbreak{}card\_classes\_filling}}}
\expandafter\newcommand\csname LN@stdClassSumcarrierXieqlineMultiplier\endcsname{\href{https://github.com/lsernau/grh-lean/blob/72fead102f663d5c7a407964a13dc0cf9d977207/grh-compact/src/New/Filling/FillingData.lean\#L26-L36}{{\leanfont Data.\allowbreak{}stdClassSum\_carrierXi\_eq\_lineMultiplier}}}
\expandafter\newcommand\csname LN@cardX\endcsname{\href{https://github.com/lsernau/grh-lean/blob/72fead102f663d5c7a407964a13dc0cf9d977207/grh-compact/src/New/Filling/FillingData.lean\#L42-L42}{{\leanfont Data.\allowbreak{}card\_X₀}}}
\expandafter\newcommand\csname LN@unitarySignFrameY\endcsname{\href{https://github.com/lsernau/grh-lean/blob/72fead102f663d5c7a407964a13dc0cf9d977207/grh-compact/src/New/Filling/FillingData.lean\#L45-L49}{{\leanfont Data.\allowbreak{}unitarySignFrameY}}}
\expandafter\newcommand\csname LN@existsbesselSystem\endcsname{\href{https://github.com/lsernau/grh-lean/blob/72fead102f663d5c7a407964a13dc0cf9d977207/grh-compact/src/New/Filling/FillingData.lean\#L56-L118}{{\leanfont Data.\allowbreak{}exists\_besselSystem}}}
\expandafter\newcommand\csname LN@instLbHPowRealNatOfNatHMulLb\endcsname{\href{https://github.com/lsernau/grh-lean/blob/72fead102f663d5c7a407964a13dc0cf9d977207/grh-compact/src/New/Filling/FillingData.lean\#L122-L123}{{\leanfont Data.\allowbreak{}instLbHPowRealNatOfNatϱHMulLb}}}
\expandafter\newcommand\csname LN@filling\endcsname{\href{https://github.com/lsernau/grh-lean/blob/72fead102f663d5c7a407964a13dc0cf9d977207/grh-compact/src/New/Filling/FillingData.lean\#L127-L140}{{\leanfont Data.\allowbreak{}filling}}}
\expandafter\newcommand\csname LN@columnLevelzeroeqcomap\endcsname{\href{https://github.com/lsernau/grh-lean/blob/72fead102f663d5c7a407964a13dc0cf9d977207/grh-compact/src/New/Filling/FillingData.lean\#L144-L149}{{\leanfont Data.\allowbreak{}columnLevel\_zero\_eq\_comap}}}
\expandafter\newcommand\csname LN@IsIntegralOff\endcsname{\href{https://github.com/lsernau/grh-lean/blob/72fead102f663d5c7a407964a13dc0cf9d977207/grh-compact/src/New/Filling/GlobalToLocal.lean\#L28-L28}{{\leanfont Filling.\allowbreak{}IsIntegralOff}}}
\expandafter\newcommand\csname LN@cosetAmplitude\endcsname{\href{https://github.com/lsernau/grh-lean/blob/72fead102f663d5c7a407964a13dc0cf9d977207/grh-compact/src/New/Filling/GlobalToLocal.lean\#L32-L34}{{\leanfont Filling.\allowbreak{}cosetAmplitude}}}
\expandafter\newcommand\csname LN@toPlaces\endcsname{\href{https://github.com/lsernau/grh-lean/blob/72fead102f663d5c7a407964a13dc0cf9d977207/grh-compact/src/New/Filling/GlobalToLocal.lean\#L41-L42}{{\leanfont Filling.\allowbreak{}toPlaces}}}
\expandafter\newcommand\csname LN@IsFillingLevel\endcsname{\href{https://github.com/lsernau/grh-lean/blob/72fead102f663d5c7a407964a13dc0cf9d977207/grh-compact/src/New/Filling/GlobalToLocal.lean\#L46-L49}{{\leanfont Filling.\allowbreak{}IsFillingLevel}}}
\expandafter\newcommand\csname LN@eventuallymemintegralOff\endcsname{\href{https://github.com/lsernau/grh-lean/blob/72fead102f663d5c7a407964a13dc0cf9d977207/grh-compact/src/New/Filling/GlobalToLocal.lean\#L52-L53}{{\leanfont Filling.\allowbreak{}eventually\_mem\_integralOff}}}
\expandafter\newcommand\csname LN@eventuallyisIntegralOff\endcsname{\href{https://github.com/lsernau/grh-lean/blob/72fead102f663d5c7a407964a13dc0cf9d977207/grh-compact/src/New/Filling/GlobalToLocal.lean\#L57-L58}{{\leanfont Filling.\allowbreak{}eventually\_isIntegralOff}}}
\expandafter\newcommand\csname LN@existsplacesintegralOff\endcsname{\href{https://github.com/lsernau/grh-lean/blob/72fead102f663d5c7a407964a13dc0cf9d977207/grh-compact/src/New/Filling/GlobalToLocal.lean\#L63-L67}{{\leanfont Filling.\allowbreak{}exists\_places\_integralOff}}}
\expandafter\newcommand\csname LN@ratToInfweylElemmul\endcsname{\href{https://github.com/lsernau/grh-lean/blob/72fead102f663d5c7a407964a13dc0cf9d977207/grh-compact/src/New/Filling/GlobalToLocal.lean\#L71-L74}{{\leanfont Filling.\allowbreak{}ratToInf\_weylElem\_mul}}}
\expandafter\newcommand\csname LN@existstoPlaceseq\endcsname{\href{https://github.com/lsernau/grh-lean/blob/72fead102f663d5c7a407964a13dc0cf9d977207/grh-compact/src/New/Filling/GlobalToLocal.lean\#L78-L89}{{\leanfont Filling.\allowbreak{}exists\_toPlaces\_eq}}}
\expandafter\newcommand\csname LN@memunitaryLeveliff\endcsname{\href{https://github.com/lsernau/grh-lean/blob/72fead102f663d5c7a407964a13dc0cf9d977207/grh-compact/src/New/Filling/GlobalToLocal.lean\#L91-L94}{{\leanfont Filling.\allowbreak{}mem\_unitaryLevel\_iff}}}
\expandafter\newcommand\csname LN@memlatStabilizeriff\endcsname{\href{https://github.com/lsernau/grh-lean/blob/72fead102f663d5c7a407964a13dc0cf9d977207/grh-compact/src/New/Filling/GlobalToLocal.lean\#L96-L97}{{\leanfont Filling.\allowbreak{}mem\_latStabilizer\_iff}}}
\expandafter\newcommand\csname LN@existsisFillingLevel\endcsname{\href{https://github.com/lsernau/grh-lean/blob/72fead102f663d5c7a407964a13dc0cf9d977207/grh-compact/src/New/Filling/GlobalToLocal.lean\#L102-L155}{{\leanfont Filling.\allowbreak{}exists\_isFillingLevel}}}
\expandafter\newcommand\csname LN@termPositionfloorGroupspec\endcsname{\href{https://github.com/lsernau/grh-lean/blob/72fead102f663d5c7a407964a13dc0cf9d977207/grh-compact/src/New/Filling/GlobalToLocal.lean\#L159-L173}{{\leanfont Filling.\allowbreak{}termPosition\_floorGroup\_spec}}}
\expandafter\newcommand\csname LN@existshasSumcosetAmplitude\endcsname{\href{https://github.com/lsernau/grh-lean/blob/72fead102f663d5c7a407964a13dc0cf9d977207/grh-compact/src/New/Filling/GlobalToLocal.lean\#L178-L212}{{\leanfont Filling.\allowbreak{}exists\_hasSum\_cosetAmplitude}}}
\expandafter\newcommand\csname LQ@existshasSumcosetAmplitude\endcsname{\leanquoteline{0}{lemma exists\_hasSum\_cosetAmplitude [WeilData E] [S.IsAdmissible] \{π : Side\} \{c s c' Λ\} (h :}\leanquoteline{2}{S.amplitudeSum (S.termsOfLink Λ) c s c' ≠ 0) (hs : ∀ γ ∈ S.termsOver c s c', ∀ l, leadingMinor}\leanquoteline{2}{(S.frameProducts γ) l ≠ 0) : ∃ (t : Eˣ) (σ : S.G E), ∃ k, (σ ∈ GT 𝔰.ϱ ∧ S.frameProducts σ 𝔰.top}\leanquoteline{2}{𝔰.top = t ∧ (∀ v, S.N \textasciicircum{} 2 < ‖Φ v t.1‖ \textasciicircum{} 2) ∧ S.reduction k = s ∧ Λ = link (weylElem 𝔰.ϱ 1 * σ)) ∧}\leanquoteline{2}{∃ L, HasSum L (S.amplitudeSum (S.termsOfLink Λ) c s c') ∧ ∀ p : Fin 𝔰.ϱ → \_, Filter.Tendsto}\leanquoteline{2}{(fun T ↦ S.cosetAmplitude T c s c' (weylElem 𝔰.ϱ (padTuple p t) * σ)) .atTop (nhds (L p))}}
\expandafter\newcommand\csname LN@tailVectorevalPlaces\endcsname{\href{https://github.com/lsernau/grh-lean/blob/72fead102f663d5c7a407964a13dc0cf9d977207/grh-compact/src/New/Filling/GlobalToLocal.lean\#L214-L217}{{\leanfont Filling.\allowbreak{}tailVector\_evalPlaces}}}
\expandafter\newcommand\csname LN@mdevalPlacesrpowhalf\endcsname{\href{https://github.com/lsernau/grh-lean/blob/72fead102f663d5c7a407964a13dc0cf9d977207/grh-compact/src/New/Filling/GlobalToLocal.lean\#L222-L229}{{\leanfont IntegralPlace.\allowbreak{}md\_evalPlaces\_rpow\_half}}}
\expandafter\newcommand\csname LN@memprincipalLeveliff\endcsname{\href{https://github.com/lsernau/grh-lean/blob/72fead102f663d5c7a407964a13dc0cf9d977207/grh-compact/src/New/Filling/GlobalToLocal.lean\#L231-L233}{{\leanfont Filling.\allowbreak{}IsFillingLevel.\allowbreak{}mem\_principalLevel\_iff}}}
\expandafter\newcommand\csname LN@reductioneqiff\endcsname{\href{https://github.com/lsernau/grh-lean/blob/72fead102f663d5c7a407964a13dc0cf9d977207/grh-compact/src/New/Filling/GlobalToLocal.lean\#L235-L238}{{\leanfont Filling.\allowbreak{}IsFillingLevel.\allowbreak{}reduction\_eq\_iff}}}
\expandafter\newcommand\csname LN@psisimpleRootSumequnipotentChar\endcsname{\href{https://github.com/lsernau/grh-lean/blob/72fead102f663d5c7a407964a13dc0cf9d977207/grh-compact/src/New/Filling/GlobalToLocal.lean\#L240-L243}{{\leanfont Filling.\allowbreak{}IsFillingLevel.\allowbreak{}psi\_simpleRootSum\_eq\_unipotentChar}}}
\expandafter\newcommand\csname LN@mdg\endcsname{\href{https://github.com/lsernau/grh-lean/blob/72fead102f663d5c7a407964a13dc0cf9d977207/grh-compact/src/New/Filling/GlobalToLocal.lean\#L245-L247}{{\leanfont Filling.\allowbreak{}IsFillingLevel.\allowbreak{}md\_g}}}
\expandafter\newcommand\csname LN@gpowmulunipotentMass\endcsname{\href{https://github.com/lsernau/grh-lean/blob/72fead102f663d5c7a407964a13dc0cf9d977207/grh-compact/src/New/Filling/GlobalToLocal.lean\#L249-L253}{{\leanfont Filling.\allowbreak{}IsFillingLevel.\allowbreak{}g\_pow\_mul\_unipotentMass}}}
\expandafter\newcommand\csname LN@cardreducedImagemul\endcsname{\href{https://github.com/lsernau/grh-lean/blob/72fead102f663d5c7a407964a13dc0cf9d977207/grh-compact/src/New/Filling/GlobalToLocal.lean\#L258-L269}{{\leanfont Filling.\allowbreak{}IsFillingLevel.\allowbreak{}card\_reducedImage\_mul}}}
\expandafter\newcommand\csname LN@cardreducedImagebesselS\endcsname{\href{https://github.com/lsernau/grh-lean/blob/72fead102f663d5c7a407964a13dc0cf9d977207/grh-compact/src/New/Filling/GlobalToLocal.lean\#L275-L296}{{\leanfont Filling.\allowbreak{}card\_reducedImage\_besselS}}}
\expandafter\newcommand\csname LN@correlationsmuleqcosetAmplitude\endcsname{\href{https://github.com/lsernau/grh-lean/blob/72fead102f663d5c7a407964a13dc0cf9d977207/grh-compact/src/New/Filling/GlobalToLocal.lean\#L302-L376}{{\leanfont Filling.\allowbreak{}IsFillingLevel.\allowbreak{}correlation\_smul\_eq\_cosetAmplitude}}}
\expandafter\newcommand\csname LQ@correlationsmuleqcosetAmplitude\endcsname{\leanquoteline{0}{lemma IsFillingLevel.correlation\_smul\_eq\_cosetAmplitude [S.IsAdmissible] [Side] \{σ p c c'\} \{Z :}\leanquoteline{2}{F.IntermediateGroup\} \{k : S.latStabilizer\} (u : σ ∈ GT 𝔰.ϱ) (o : S.IsIntegralOff T}\leanquoteline{2}{(S.termPosition c c' (weylElem 𝔰.ϱ p * σ))) \{a b : Z.Cosets\} (e : (S.toPlaces T k)⁻¹ * a.out.1 *}\leanquoteline{2}{b.out.1⁻¹ ∈ F.congGroup) : let q := Units.map (ratToPlaces T) ∘ p; let γ := weylElem 𝔰.ϱ p * σ}\leanquoteline{4}{F.correlation Z q (S.toPlaces T (S.termPosition c c' σ))⁻¹ (a, b) • besselCorrelation volume}\leanquoteline{6}{(.toLp 2 <| S.besselFn ·) (S.ratToInf γ) = (F.headDensity 𝔰.ϱ q 1).toReal •}\leanquoteline{6}{F.groupMeasure.real F.congGroup • S.cosetAmplitude T c (S.reduction k) c' γ}}
\expandafter\newcommand\csname LN@instAdmissible\endcsname{\href{https://github.com/lsernau/grh-lean/blob/72fead102f663d5c7a407964a13dc0cf9d977207/grh-compact/src/New/Filling/FillingAdmissible.lean\#L40-L81}{{\leanfont instance : D.\allowbreak{}filling.\allowbreak{}IsAdmissible}}}
\newcommand{\ConcordanceRows}{\crow{relIndexlatStabilizer}{\href{https://github.com/lsernau/grh-lean/blob/72fead102f663d5c7a407964a13dc0cf9d977207/grh-compact/src/New/Filling/FillingAdmissible.lean\#L16-L28}{{\leanfont Data.\allowbreak{}relIndex\_latStabilizer}}}{FillingAdmissible.lean}{16--28}
\crow{cardclassesfilling}{\href{https://github.com/lsernau/grh-lean/blob/72fead102f663d5c7a407964a13dc0cf9d977207/grh-compact/src/New/Filling/FillingAdmissible.lean\#L30-L37}{{\leanfont Data.\allowbreak{}card\_classes\_filling}}}{FillingAdmissible.lean}{30--37}
\crow{stdClassSumcarrierXieqlineMultiplier}{\href{https://github.com/lsernau/grh-lean/blob/72fead102f663d5c7a407964a13dc0cf9d977207/grh-compact/src/New/Filling/FillingData.lean\#L26-L36}{{\leanfont Data.\allowbreak{}stdClassSum\_carrierXi\_eq\_lineMultiplier}}}{FillingData.lean}{26--36}
\crow{cardX}{\href{https://github.com/lsernau/grh-lean/blob/72fead102f663d5c7a407964a13dc0cf9d977207/grh-compact/src/New/Filling/FillingData.lean\#L42-L42}{{\leanfont Data.\allowbreak{}card\_X₀}}}{FillingData.lean}{42--42}
\crow{unitarySignFrameY}{\href{https://github.com/lsernau/grh-lean/blob/72fead102f663d5c7a407964a13dc0cf9d977207/grh-compact/src/New/Filling/FillingData.lean\#L45-L49}{{\leanfont Data.\allowbreak{}unitarySignFrameY}}}{FillingData.lean}{45--49}
\crow{existsbesselSystem}{\href{https://github.com/lsernau/grh-lean/blob/72fead102f663d5c7a407964a13dc0cf9d977207/grh-compact/src/New/Filling/FillingData.lean\#L56-L118}{{\leanfont Data.\allowbreak{}exists\_besselSystem}}}{FillingData.lean}{56--118}
\crow{instLbHPowRealNatOfNatHMulLb}{\href{https://github.com/lsernau/grh-lean/blob/72fead102f663d5c7a407964a13dc0cf9d977207/grh-compact/src/New/Filling/FillingData.lean\#L122-L123}{{\leanfont Data.\allowbreak{}instLbHPowRealNatOfNatϱHMulLb}}}{FillingData.lean}{122--123}
\crow{filling}{\href{https://github.com/lsernau/grh-lean/blob/72fead102f663d5c7a407964a13dc0cf9d977207/grh-compact/src/New/Filling/FillingData.lean\#L127-L140}{{\leanfont Data.\allowbreak{}filling}}}{FillingData.lean}{127--140}
\crow{columnLevelzeroeqcomap}{\href{https://github.com/lsernau/grh-lean/blob/72fead102f663d5c7a407964a13dc0cf9d977207/grh-compact/src/New/Filling/FillingData.lean\#L144-L149}{{\leanfont Data.\allowbreak{}columnLevel\_zero\_eq\_comap}}}{FillingData.lean}{144--149}
\crow{IsIntegralOff}{\href{https://github.com/lsernau/grh-lean/blob/72fead102f663d5c7a407964a13dc0cf9d977207/grh-compact/src/New/Filling/GlobalToLocal.lean\#L28-L28}{{\leanfont Filling.\allowbreak{}IsIntegralOff}}}{GlobalToLocal.lean}{28--28}
\crow{cosetAmplitude}{\href{https://github.com/lsernau/grh-lean/blob/72fead102f663d5c7a407964a13dc0cf9d977207/grh-compact/src/New/Filling/GlobalToLocal.lean\#L32-L34}{{\leanfont Filling.\allowbreak{}cosetAmplitude}}}{GlobalToLocal.lean}{32--34}
\crow{toPlaces}{\href{https://github.com/lsernau/grh-lean/blob/72fead102f663d5c7a407964a13dc0cf9d977207/grh-compact/src/New/Filling/GlobalToLocal.lean\#L41-L42}{{\leanfont Filling.\allowbreak{}toPlaces}}}{GlobalToLocal.lean}{41--42}
\crow{IsFillingLevel}{\href{https://github.com/lsernau/grh-lean/blob/72fead102f663d5c7a407964a13dc0cf9d977207/grh-compact/src/New/Filling/GlobalToLocal.lean\#L46-L49}{{\leanfont Filling.\allowbreak{}IsFillingLevel}}}{GlobalToLocal.lean}{46--49}
\crow{eventuallymemintegralOff}{\href{https://github.com/lsernau/grh-lean/blob/72fead102f663d5c7a407964a13dc0cf9d977207/grh-compact/src/New/Filling/GlobalToLocal.lean\#L52-L53}{{\leanfont Filling.\allowbreak{}eventually\_mem\_integralOff}}}{GlobalToLocal.lean}{52--53}
\crow{eventuallyisIntegralOff}{\href{https://github.com/lsernau/grh-lean/blob/72fead102f663d5c7a407964a13dc0cf9d977207/grh-compact/src/New/Filling/GlobalToLocal.lean\#L57-L58}{{\leanfont Filling.\allowbreak{}eventually\_isIntegralOff}}}{GlobalToLocal.lean}{57--58}
\crow{existsplacesintegralOff}{\href{https://github.com/lsernau/grh-lean/blob/72fead102f663d5c7a407964a13dc0cf9d977207/grh-compact/src/New/Filling/GlobalToLocal.lean\#L63-L67}{{\leanfont Filling.\allowbreak{}exists\_places\_integralOff}}}{GlobalToLocal.lean}{63--67}
\crow{ratToInfweylElemmul}{\href{https://github.com/lsernau/grh-lean/blob/72fead102f663d5c7a407964a13dc0cf9d977207/grh-compact/src/New/Filling/GlobalToLocal.lean\#L71-L74}{{\leanfont Filling.\allowbreak{}ratToInf\_weylElem\_mul}}}{GlobalToLocal.lean}{71--74}
\crow{existstoPlaceseq}{\href{https://github.com/lsernau/grh-lean/blob/72fead102f663d5c7a407964a13dc0cf9d977207/grh-compact/src/New/Filling/GlobalToLocal.lean\#L78-L89}{{\leanfont Filling.\allowbreak{}exists\_toPlaces\_eq}}}{GlobalToLocal.lean}{78--89}
\crow{memunitaryLeveliff}{\href{https://github.com/lsernau/grh-lean/blob/72fead102f663d5c7a407964a13dc0cf9d977207/grh-compact/src/New/Filling/GlobalToLocal.lean\#L91-L94}{{\leanfont Filling.\allowbreak{}mem\_unitaryLevel\_iff}}}{GlobalToLocal.lean}{91--94}
\crow{memlatStabilizeriff}{\href{https://github.com/lsernau/grh-lean/blob/72fead102f663d5c7a407964a13dc0cf9d977207/grh-compact/src/New/Filling/GlobalToLocal.lean\#L96-L97}{{\leanfont Filling.\allowbreak{}mem\_latStabilizer\_iff}}}{GlobalToLocal.lean}{96--97}
\crow{existsisFillingLevel}{\href{https://github.com/lsernau/grh-lean/blob/72fead102f663d5c7a407964a13dc0cf9d977207/grh-compact/src/New/Filling/GlobalToLocal.lean\#L102-L155}{{\leanfont Filling.\allowbreak{}exists\_isFillingLevel}}}{GlobalToLocal.lean}{102--155}
\crow{termPositionfloorGroupspec}{\href{https://github.com/lsernau/grh-lean/blob/72fead102f663d5c7a407964a13dc0cf9d977207/grh-compact/src/New/Filling/GlobalToLocal.lean\#L159-L173}{{\leanfont Filling.\allowbreak{}termPosition\_floorGroup\_spec}}}{GlobalToLocal.lean}{159--173}
\crow{existshasSumcosetAmplitude}{\href{https://github.com/lsernau/grh-lean/blob/72fead102f663d5c7a407964a13dc0cf9d977207/grh-compact/src/New/Filling/GlobalToLocal.lean\#L178-L212}{{\leanfont Filling.\allowbreak{}exists\_hasSum\_cosetAmplitude}}}{GlobalToLocal.lean}{178--212}
\crow{tailVectorevalPlaces}{\href{https://github.com/lsernau/grh-lean/blob/72fead102f663d5c7a407964a13dc0cf9d977207/grh-compact/src/New/Filling/GlobalToLocal.lean\#L214-L217}{{\leanfont Filling.\allowbreak{}tailVector\_evalPlaces}}}{GlobalToLocal.lean}{214--217}
\crow{mdevalPlacesrpowhalf}{\href{https://github.com/lsernau/grh-lean/blob/72fead102f663d5c7a407964a13dc0cf9d977207/grh-compact/src/New/Filling/GlobalToLocal.lean\#L222-L229}{{\leanfont IntegralPlace.\allowbreak{}md\_evalPlaces\_rpow\_half}}}{GlobalToLocal.lean}{222--229}
\crow{memprincipalLeveliff}{\href{https://github.com/lsernau/grh-lean/blob/72fead102f663d5c7a407964a13dc0cf9d977207/grh-compact/src/New/Filling/GlobalToLocal.lean\#L231-L233}{{\leanfont Filling.\allowbreak{}IsFillingLevel.\allowbreak{}mem\_principalLevel\_iff}}}{GlobalToLocal.lean}{231--233}
\crow{reductioneqiff}{\href{https://github.com/lsernau/grh-lean/blob/72fead102f663d5c7a407964a13dc0cf9d977207/grh-compact/src/New/Filling/GlobalToLocal.lean\#L235-L238}{{\leanfont Filling.\allowbreak{}IsFillingLevel.\allowbreak{}reduction\_eq\_iff}}}{GlobalToLocal.lean}{235--238}
\crow{psisimpleRootSumequnipotentChar}{\href{https://github.com/lsernau/grh-lean/blob/72fead102f663d5c7a407964a13dc0cf9d977207/grh-compact/src/New/Filling/GlobalToLocal.lean\#L240-L243}{{\leanfont Filling.\allowbreak{}IsFillingLevel.\allowbreak{}psi\_simpleRootSum\_eq\_unipotentChar}}}{GlobalToLocal.lean}{240--243}
\crow{mdg}{\href{https://github.com/lsernau/grh-lean/blob/72fead102f663d5c7a407964a13dc0cf9d977207/grh-compact/src/New/Filling/GlobalToLocal.lean\#L245-L247}{{\leanfont Filling.\allowbreak{}IsFillingLevel.\allowbreak{}md\_g}}}{GlobalToLocal.lean}{245--247}
\crow{gpowmulunipotentMass}{\href{https://github.com/lsernau/grh-lean/blob/72fead102f663d5c7a407964a13dc0cf9d977207/grh-compact/src/New/Filling/GlobalToLocal.lean\#L249-L253}{{\leanfont Filling.\allowbreak{}IsFillingLevel.\allowbreak{}g\_pow\_mul\_unipotentMass}}}{GlobalToLocal.lean}{249--253}
\crow{cardreducedImagemul}{\href{https://github.com/lsernau/grh-lean/blob/72fead102f663d5c7a407964a13dc0cf9d977207/grh-compact/src/New/Filling/GlobalToLocal.lean\#L258-L269}{{\leanfont Filling.\allowbreak{}IsFillingLevel.\allowbreak{}card\_reducedImage\_mul}}}{GlobalToLocal.lean}{258--269}
\crow{cardreducedImagebesselS}{\href{https://github.com/lsernau/grh-lean/blob/72fead102f663d5c7a407964a13dc0cf9d977207/grh-compact/src/New/Filling/GlobalToLocal.lean\#L275-L296}{{\leanfont Filling.\allowbreak{}card\_reducedImage\_besselS}}}{GlobalToLocal.lean}{275--296}
\crow{correlationsmuleqcosetAmplitude}{\href{https://github.com/lsernau/grh-lean/blob/72fead102f663d5c7a407964a13dc0cf9d977207/grh-compact/src/New/Filling/GlobalToLocal.lean\#L302-L376}{{\leanfont Filling.\allowbreak{}IsFillingLevel.\allowbreak{}correlation\_smul\_eq\_cosetAmplitude}}}{GlobalToLocal.lean}{302--376}
\crow{instAdmissible}{\href{https://github.com/lsernau/grh-lean/blob/72fead102f663d5c7a407964a13dc0cf9d977207/grh-compact/src/New/Filling/FillingAdmissible.lean\#L40-L81}{{\leanfont instance : D.\allowbreak{}filling.\allowbreak{}IsAdmissible}}}{FillingAdmissible.lean}{40--81}}
% the notations of the defined functions: one slot for each explicit argument of the Lean declaration
\newcommand{\SD}[1]{S(#1)}
\newcommand{\intoff}[3]{#3\in\operatorname{Int}^{#2}(#1)}
\newcommand{\cosamp}[6]{a^{\sigma}_{#1,#2}(#3,#4,#5;#6)}
\newcommand{\loc}[2]{\operatorname{loc}_{#1,#2}}
\newcommand{\isFL}[3]{#3\in\operatorname{FL}(#1,#2)}
\newcommand{\UnitMembers}{497}
\newcommand{\UnitInNote}{34}
\newcommand{\UnitProofs}{179}
\newcommand{\UnitSteps}{1371}
\newcommand{\UnitAllSteps}{1479}


\linespread{1.08}                       % 11 pt on 14.7 pt
\setlength{\parindent}{0pt}
\setlength{\parskip}{5.5pt plus 1.5pt minus 0.5pt}
\setlength{\abovedisplayskip}{7.5pt plus 2pt minus 3pt}
\setlength{\belowdisplayskip}{7.5pt plus 2pt minus 3pt}
\setlength{\abovedisplayshortskip}{3pt plus 2pt}
\setlength{\belowdisplayshortskip}{6pt plus 2pt minus 2pt}
\emergencystretch=1.5em
\clubpenalty=10000 \widowpenalty=10000 \displaywidowpenalty=10000
\relpenalty=7000 \binoppenalty=9000
\raggedbottom
\pagestyle{fancy}\fancyhf{}\renewcommand{\headrulewidth}{0pt}
\fancyfoot[C]{\sffamily\fontsize{7.6}{9}\selectfont\color{inkfaint}\thepage\,/\,\pageref*{LastPage}}
\setlist{topsep=2pt,partopsep=0pt,parsep=0pt,itemsep=4pt,leftmargin=2.1em,labelsep=0.55em}

\newcommand{\Q}{\mathbb{Q}}
\newcommand{\Z}{\mathbb{Z}}
\newcommand{\F}{\mathbb{F}}
\newcommand{\len}{\operatorname{len}}
\newcommand{\cok}{\operatorname{cok}}
\newcommand{\End}{\operatorname{End}}
\newcommand{\Hom}{\operatorname{Hom}}
\newcommand{\Gal}{\operatorname{Gal}}
\newcommand{\diag}{\operatorname{diag}}
\newcommand{\kw}[1]{\emph{#1}}
\newcommand{\m}[1]{\mbox{$#1$}}        % a short formula that must not be broken across lines
\newcommand{\head}[1]{\par\addvspace{17pt plus 4pt minus 2pt}\Needspace{6\baselineskip}%
  \noindent{\fontsize{13.2}{16}\selectfont\bfseries\color{ink}#1}\par\nopagebreak\vspace{-1pt}}
\newcommand{\sect}[2]{\refstepcounter{section}\label{#2}\par\addvspace{19pt plus 4pt minus 2pt}\Needspace{7\baselineskip}%
  \noindent\llap{{\fontsize{13.2}{16}\selectfont\bfseries\color{inkfaint}\thesection}\hspace{9pt}}%
  {\fontsize{13.2}{16}\selectfont\bfseries\color{ink}#1}\par\nopagebreak\vspace{-1pt}}
\newcommand{\subhead}[1]{\par\addvspace{9pt plus 2pt}\Needspace{4\baselineskip}\noindent\textbf{#1}\enspace\ignorespaces}
% where a statement lives in the Lean archive
\newcommand{\leantag}{{\sffamily\bfseries\fontsize{6.3}{8}\selectfont\color{leancol}\addfontfeatures{LetterSpace=7}LEAN}}
\newcommand{\sep}{\unskip\nobreak\hspace{5pt}{\color{inkfaint}\textperiodcentered}\hspace{5pt plus 2pt}\ignorespaces}
\newcommand{\leanline}[1]{\par\nopagebreak\vspace{-2.5pt}{\fontsize{8.9}{12.4}\selectfont\color{inksoft}\raggedright
  \noindent\llap{\leantag\hspace{9pt}}#1\par}\vspace{1pt}}
% a displayed statement in a tinted box: #1 colour, #2 text
\newcommand{\statement}[2]{\par\addvspace{7pt}\begingroup\setlength{\fboxsep}{8.5pt}\setlength{\fboxrule}{0.6pt}%
  \noindent\fcolorbox{#1!30!white}{#1!3!white}{\parbox{\dimexpr\linewidth-2\fboxsep-2\fboxrule\relax}{#2}}%
  \endgroup\par\addvspace{7pt}}
\newcommand{\nt}[2]{\par\noindent\hangindent=1.1em\hangafter=1 #1\enspace{\color{inksoft}#2}\par\vspace{1pt}}
\newcommand{\legdot}[1]{\tikz[baseline=-0.62ex]\fill[#1] (0,0) circle (2.5pt);\hspace{3pt}}
\newcommand{\crow}[4]{\csname CS@#1\endcsname & #2 & \textcolor{inksoft}{#3} & \textcolor{inksoft}{#4}\\}
\newcommand{\leanquoteline}[2]{\par\noindent\hangindent=\dimexpr#1\fontdimen2\font+2.2em\relax\hangafter=1
  \hspace*{#1\fontdimen2\font}#2}
\newcommand{\quoted}[1]{\par\vspace{3pt}{\ttfamily\fontsize{7.5}{10.4}\selectfont\color{leancol}\raggedright
  \setlength{\parskip}{0pt}\leanquote{#1}\par}\vspace{4pt}}
\newcommand{\bx}[3]{{\sffamily\bfseries\fontsize{7.3}{9.2}\selectfont\color{#1!82!black}#2}\\[1.5pt]#3}
\newcommand{\cs}[2]{\expandafter\newcommand\csname CS@#1\endcsname{#2}}

% theorems
\newtheoremstyle{plainx}{7pt plus 2pt}{7pt plus 2pt}{\itshape}{}{\bfseries}{.}{0.5em}{}
\newtheoremstyle{uprightx}{7pt plus 2pt}{7pt plus 2pt}{}{}{\bfseries}{.}{0.5em}{}
\theoremstyle{plainx}
\newtheorem{theorem}{Theorem}[section]
\newtheorem{lemma}[theorem]{Lemma}
\newtheorem{proposition}[theorem]{Proposition}
\newtheorem{corollary}[theorem]{Corollary}
\theoremstyle{uprightx}
\newtheorem{definition}[theorem]{Definition}
\newtheorem{setting}[theorem]{Setting}          % hypotheses shared by several statements; each of them cites it by number
\newtheorem{remark}[theorem]{Remark}
\newtheorem{example}[theorem]{Example}
\renewcommand{\qedsymbol}{{\setlength{\fboxsep}{0pt}\setlength{\fboxrule}{0.5pt}\fbox{\rule{0pt}{0.46em}\hspace{0.46em}}}}
% what holds throughout the note: the ambient objects only (implicit and instance arguments), never a hypothesis
\newenvironment{conventions}{\par\addvspace{5pt}\noindent\textbf{Conventions.}\enspace\ignorespaces}{\par\addvspace{3pt}}
\makeatletter
\renewenvironment{proof}[1][Proof]{\par\pushQED{\qed}\normalfont\topsep0pt\partopsep0pt
  \trivlist\item[\hskip\labelsep\itshape #1.]\ignorespaces}{\popQED\endtrivlist\@endpefalse}
\makeatother

% where each declaration of the concordance is treated: \cs{key}{place}, the place naming a \label or a lettered item
\cs{IsIntegralOff}{Definition~\ref{d:off}}
\cs{toPlaces}{Definition~\ref{d:off}}
\cs{IsFillingLevel}{Definition~\ref{d:FL}}
\cs{cosetAmplitude}{Definition~\ref{d:cosamp}}
\cs{eventuallymemintegralOff}{Theorem C}
\cs{eventuallyisIntegralOff}{Theorem C}
\cs{existsplacesintegralOff}{Theorem C}
\cs{existsisFillingLevel}{Theorem C}
\cs{existstoPlaceseq}{Lemma~\ref{l:lift}}
\cs{memunitaryLeveliff}{Lemma~\ref{l:levels}}
\cs{memlatStabilizeriff}{Lemma~\ref{l:levels}}
\cs{memprincipalLeveliff}{Lemma~\ref{l:levels}}
\cs{reductioneqiff}{Lemma~\ref{l:levels}}
\cs{psisimpleRootSumequnipotentChar}{Lemma~\ref{l:char}}
\cs{mdevalPlacesrpowhalf}{Lemma~\ref{l:mod}}
\cs{mdg}{Lemma~\ref{l:mod}}
\cs{gpowmulunipotentMass}{Lemma~\ref{l:mod}}
\cs{cardreducedImagemul}{Proposition~\ref{p:count}}
\cs{cardreducedImagebesselS}{Corollary~\ref{c:bessel}}
\cs{termPositionfloorGroupspec}{Lemma~\ref{l:floor}}
\cs{ratToInfweylElemmul}{Lemma~\ref{l:weylinf}}
\cs{tailVectorevalPlaces}{Lemma~\ref{l:weylinf}}
\cs{existshasSumcosetAmplitude}{Theorem B}
\cs{correlationsmuleqcosetAmplitude}{Theorem A}
\cs{filling}{Definition~\ref{d:SD}}
\cs{cardX}{Lemma~\ref{l:small}}
\cs{unitarySignFrameY}{Lemma~\ref{l:small}}
\cs{instLbHPowRealNatOfNatHMulLb}{Lemma~\ref{l:small}}
\cs{existsbesselSystem}{Proposition~\ref{p:bessel}}
\cs{columnLevelzeroeqcomap}{Theorem D}
\cs{relIndexlatStabilizer}{Theorem D}
\cs{cardclassesfilling}{Theorem D}
\cs{stdClassSumcarrierXieqlineMultiplier}{Lemma~\ref{l:poly}}
\cs{instAdmissible}{Theorem D}
% ----- macros-148: the recalled notions of unit 148 (chunks a, b) and of its neighbours, with every explicit
% ----- argument of the Lean definition as a slot. Shared by Notes F, P, T, G; copy unchanged into the preamble.
\newcommand{\fs}{\mathfrak{s}}                      % the shape (rho, kappa, k_0)
\newcommand{\sT}{\mathsf{T}}                        % a finite set of finite places
\newcommand{\cL}{\mathcal{L}}                       % a level
\newcommand{\AEf}{\mathbb{A}_{E,f}}
\newcommand{\hO}{\hat{O}_E}
\newcommand{\GL}{\mathrm{GL}}
\newcommand{\full}{\mathrm{full}}
\newcommand{\cut}{\mathrm{cut}}
\newcommand{\rh}[1]{\varrho(#1)}                    % Shape.ϱ
\newcommand{\kk}[1]{k(#1)}                          % Shape.k = #κ
\newcommand{\nn}[1]{n(#1)}                          % Shape.n = ϱ + 1
\newcommand{\rr}[1]{r(#1)}                          % Shape.r = 2ϱ + k + 1
\newcommand{\Ix}[1]{I(#1)}                          % Shape.BesselIndex
\newcommand{\Mx}[1]{M(#1)}                          % Shape.MiddleIndex
\newcommand{\ninf}[1]{f(#1)}                        % numInfPlaces E
\newcommand{\ET}[1]{E_{#1}}                         % LocProd T
% the components of the filling data S
\newcommand{\bS}[1]{b(#1)}                          % S.b
\newcommand{\NS}[1]{N(#1)}                          % S.N
\newcommand{\qS}[2]{q_{#2}(#1)}                     % S.levelPrime i
\newcommand{\gS}[1]{g(#1)}                          % S.g = q_0 q_1
\newcommand{\piS}[1]{\bar\pi(#1)}                   % S.piBar
\newcommand{\vpS}[1]{\varpi(#1)}                    % S.ϖ
\newcommand{\ChS}[2]{\mathrm{Ch}_{#2}(#1)}          % S.chosenPlaces i
\newcommand{\KcolS}[1]{K^{\mathrm{col}}(#1)}        % S.columnLevel
\newcommand{\BsysS}[2]{\mathfrak{B}_{#2}(#1)}       % S.besselSys v
% forms, groups, levels of the set-up (Setup.lean)
\newcommand{\Trat}[2]{T^{\mathrm{rat}}_{#1}(#2)}    % S.ratMiddle R
\newcommand{\Tlat}[2]{T^{\mathrm{lat}}_{#1}(#2)}    % S.latMiddle R
\newcommand{\GS}[2]{G_{#1}(#2)}                     % S.G R
\newcommand{\GFin}[1]{G^{\mathrm{lat}}_{#1}}        % S.GFin
\newcommand{\GRes}[1]{\mathsf{G}(#1)}               % S.GRes
\newcommand{\Cl}[1]{\mathrm{Cl}(#1)}                % S.Classes
\newcommand{\Bc}[1]{\mathrm{Bc}(#1)}                % S.BesselCoord
\newcommand{\KS}[1]{K(#1)}                          % S.latStabilizer
\newcommand{\redS}[1]{\mathrm{red}_{#1}}            % S.reduction
\newcommand{\KgS}[1]{K_g(#1)}                       % S.principalLevel
\newcommand{\redim}[2]{\overline{#2}^{\,#1}}        % S.reducedImage A
\newcommand{\kapS}[2]{\kappa_{#2}(#1)}              % S.kappa v
\newcommand{\gid}[1]{\hat{g}(#1)}                   % S.gIdele
\newcommand{\latB}[1]{\Lambda(#1)}                  % S.latBasis
\newcommand{\rinf}[1]{\mathrm{r}^{\infty}_{#1}}     % S.ratToInf
\newcommand{\rfin}[1]{\mathrm{r}^{f}_{#1}}          % S.ratToFin
\newcommand{\incH}[1]{\mathrm{inc}_{#1}}            % S.inclHFin
\newcommand{\hc}[2]{h_{#1}(#2)}                     % S.h c
\newcommand{\pos}[4]{\mathrm{pos}_{#1}(#2,#3;#4)}   % S.termPosition c c' γ
\newcommand{\terms}[4]{\mathrm{Tm}_{#1}(#2,#3,#4)}  % S.termsOver c s c'
\newcommand{\Zf}[2]{Z_{#1}(#2)}                     % S.frameProducts γ
\newcommand{\tdisc}[2]{\mathsf{g}_{#1}(#2)}         % S.termDisc γ
\newcommand{\clU}[2]{\mathcal{U}_{#1,#2}}           % S.classUnip c
\newcommand{\cutoff}[2]{\beta_{#1,#2}}              % S.classCutoff c
\newcommand{\besfn}[4]{W^{#1}_{#2}(#3)_{#4}}        % S.besselFn y i, with the side
\newcommand{\amp}[5]{a^{#1}_{#2}(#3,#4;#5)}         % S.amplitude c c' γ, with the side
\newcommand{\ampabs}[5]{|a|^{#1}_{#2}(#3,#4;#5)}    % S.amplitudeAbs c c' γ, with the side
\newcommand{\ampsum}[6]{A^{#1}_{#2}(#3;#4,#5,#6)}   % S.amplitudeSum T c s c', with the side
\newcommand{\uchar}[2]{\psi^{\mathcal{U}}_{#1}(#2)} % S.unipotentChar u
\newcommand{\ucov}[1]{\nu(#1)}                      % S.unipCovolume
\newcommand{\siegC}[1]{c_{*}(#1)}                   % S.siegelConst
\newcommand{\packC}[1]{a_{*}(#1)}                   % S.packetConst
\newcommand{\uord}[1]{\mathrm{uo}(#1)}              % S.unipOrder
\newcommand{\cdim}[1]{\mathrm{cd}(#1)}              % S.columnDim
\newcommand{\bdeg}[1]{\deg(#1)}                     % S.besselDegree
\newcommand{\mainE}[1]{\mathsf{d}(#1)}              % S.mainEnergy
\newcommand{\besE}[1]{\mathsf{d}_{1}(#1)}           % S.besselEnergy
\newcommand{\tailS}[1]{\mathrm{tl}(#1)}             % S.tailScale
\newcommand{\codimF}[1]{\mathrm{cf}(#1)}            % S.codimFactor
\newcommand{\fscale}[3]{\mathrm{fs}_{#1}(#2,#3)}    % S.frameScaleSq v a
\newcommand{\fsprod}[1]{\mathrm{fs}(#1)}            % S.frameScaleProd
\newcommand{\mass}[2]{\operatorname{mass}(#1,#2)}   % mass E r
% chunk b and the sieve
\newcommand{\energy}[3]{\mathrm{En}_{#1}(#2;#3)}    % Filling.energy A Ω (S implicit in Lean, shown)
\newcommand{\linkE}[5]{\mathrm{LE}^{#1}_{#2}(#3;#4,#5)} % Filling.linkEnergy T Ω' R (side and S implicit, shown)
\newcommand{\fixset}[3]{\mathrm{Fix}_{#1}(#2,#3)}   % Filling.fixingSet Sg d (S implicit, shown)
\newcommand{\vfix}[3]{\mathrm{Fix}_{#1,#2}(#3)}     % Filling.vectorFixer w d (S implicit, shown)
\newcommand{\charT}[2]{\psi_{#1,#2}}                % S.charOnPlaces T
\newcommand{\fgain}[3]{\beta_{#1,#2}(#3)}           % S.frequencyGain T ξ
\newcommand{\link}[1]{\operatorname{link}(#1)}      % link γ
\newcommand{\clink}[2]{\operatorname{link}^{\circ}(#1,#2)} % commonLink γ γ'
\newcommand{\tdnorm}[2]{\|\mathsf{g}\|_{#1}(#2)}    % S.termDiscNorm γ
\newcommand{\mnorm}[3]{\mathrm{mn}_{#1}(#2;#3)}     % S.minorNorm j γ
\newcommand{\cgauge}[5]{\mathrm{cg}_{#1}(#2;#3,#4;#5)} % S.classGauge j c c' γ
\newcommand{\pinv}[5]{\mathrm{pi}_{#1}(#2,#3;#4,#5)} % S.pairInvariant c c' γ γ'
\newcommand{\ldsq}[3]{\mathrm{ld}_{#1}(#2,#3)}      % S.linkDiscSq v e
\newcommand{\lmet}[1]{Q^{\mathrm{link}}_{#1}}       % S.linkMetric
% ----- end of macros-148
% ----- notions of the background library, of the set-up of this unit, of Note T and of other units that this
% ----- note recalls, written by hand with every explicit argument (TeX as in the notes that state them)
\newcommand{\JT}[1]{J(#1)}                          % Shape.Form.splitGram
\newcommand{\GT}[1]{G(#1)}                          % Shape.Form.G
\newcommand{\UT}[1]{\mathcal{U}(#1)}                % besselU
\newcommand{\HT}[1]{\mathcal{H}(#1)}                % besselH
\newcommand{\ST}[1]{\mathcal{S}(#1)}                % besselS
\newcommand{\Hl}[2]{\mathcal{H}_{#1}(#2)}           % GT l
\newcommand{\sd}[1]{\operatorname{sd}(#1)}          % simpleRootSum
\newcommand{\ZT}[2]{Z_{#1}(#2)}                     % Shape.Form.frameGram
\newcommand{\weyl}[2]{w_{#1}(#2)}                   % weylElem l p
\newcommand{\pad}[2]{(#1\,|\,#2)}                   % padTuple p t
\newcommand{\tailv}[3]{\operatorname{tv}(#1,#2,#3)} % tailVector π t x
\newcommand{\gix}[1]{[#1]}                          % genIndex
\newcommand{\modT}[1]{\operatorname{mod}_{#1}}      % the modulus of the place structure on E_T
\newcommand{\psiT}[1]{\psi_{#1}}                    % the character of the place structure on E_T
\newcommand{\KA}[2]{\hat{K}_{#1}(#2)}               % unitaryLevel (integral adeles) u (split Gram matrix of latMiddle over the finite adeles)
\newcommand{\KT}[3]{K_{#1,#2}(#3)}                  % unitaryLevel (integers of E_T) u (split Gram matrix of latMiddle over E_T)
\newcommand{\cgp}[1]{\mathrm{K}^{\circ}(#1)}        % Level.congGroup F
\newcommand{\gL}[1]{g_{#1}}                         % the unit of a level
\newcommand{\IsStd}[1]{#1\in\operatorname{StdLev}}  % IntegralPlace.IsStdLevel F
\newcommand{\Levels}[3]{\operatorname{Lev}(#1,#2,#3)} % Level ι T m
\newcommand{\gM}[1]{\mathsf{M}(#1)}                 % Level.groupMeasure F
\newcommand{\etab}[1]{\eta_{<#1}}                   % ascConv heisenbergMeasure j
\newcommand{\fmod}[3]{\mathsf{d}_{#2}(#1;#3)}       % Level.floorModulus F l g
\newcommand{\hdens}[4]{\mathsf{h}_{#2}(#1;#3,#4)}   % Level.headDensity F j p q
\newcommand{\bcorr}[3]{\mathrm{J}(#1,#2;#3)}        % besselCorrelation μ W γ
\newcommand{\phs}[2]{\chi_{#1}(#2)}                 % IntegralPlace.phase g z (Note T)
\newcommand{\IntG}[1]{\operatorname{Int}(#1)}       % Level.IntermediateGroup F (Note T)
\newcommand{\umass}[1]{\mathsf{v}(#1)}              % Level.unipotentMass F (Note T)
\newcommand{\Cos}[1]{\operatorname{Cos}(#1)}        % Level.IntermediateGroup.Cosets Z (Note T)
\newcommand{\bvec}[2]{W_{#1}(#2)}                   % Level.IntermediateGroup.besselVector Z y (Note T)
\newcommand{\corr}[3]{\mathcal{J}(#1;#2,#3)}        % Level.correlation Z p l (Note T)
\newcommand{\tlink}[2]{\mathrm{Tm}^{\mathrm{lk}}_{#1}(#2)} % Filling.termsOfLink S Λ
\newcommand{\bker}[4]{\mathrm{bk}^{#1}_{#2}(#3,#4)} % Filling.besselKernel γ y, with the side and S
\newcommand{\Fits}[6]{\operatorname{Fits}(#1;#2,#3,#4,#5,#6)} % Filling.Fits S V l kw Y η
\newcommand{\sGD}[1]{\mathsf{G}(#1)}                % Data.G
\newcommand{\Colv}[2]{K^{\mathrm{col}}_{#2}(#1)}    % Data.columnLevel D j
\newcommand{\IsBes}[3]{#1\in\operatorname{Bes}(#2,#3)} % Filling.IsBessel B l kw
\newcommand{\basisVec}[2]{\operatorname{bv}_{#1}(#2)} % Data.basisVec D i
\newcommand{\posEquiv}[1]{\operatorname{pe}_D(#1)}  % Data.posEquiv j
\newcommand{\detRows}[4]{\operatorname{dr}_D(#1,#2,#3,#4)} % Data.detRows j i Y τ
\newcommand{\linePolyCoords}[3]{\operatorname{lc}_D(#1,#2,#3)} % Data.linePolyCoords j i Y
\newcommand{\psiS}[1]{\psi_{#1}}                     % Filling.selfDualChar S
\newcommand{\vol}{\operatorname{vol}}
\newcommand{\rePart}{\operatorname{Re}}
\newcommand{\CC}{\mathbb{C}}
\newcommand{\RR}{\mathbb{R}}
\newcommand{\NN}{\mathbb{N}}
\newcommand{\Inp}[1]{Input~\ref{i:#1}}

\begin{document}
\color{ink}

{\raggedright
{\sffamily\bfseries\fontsize{7.2}{9}\selectfont\color{leancol}\addfontfeatures{LetterSpace=8}READBACK NOTE 38\par}
\vspace{5pt}
{\fontsize{21.5}{25.5}\selectfont\bfseries From amplitude sums to local correlations: filling levels at finitely many places, and the filling data of the admissible data\par}
\vspace{5pt}
{\fontsize{12.4}{16}\selectfont\itshape\color{inksoft} One unit of the Lean artifact \ArtifactName, read back into
mathematics\par}}
\vspace{8pt}
{\color{rulegrey}\hrule height 0.6pt}
\vspace{7pt}
\thispagestyle{empty}

\textbf{Summary.} The second moment of the Poincaré series is a sum of \emph{amplitudes} of rational elements
of a unitary group, sorted by classes and by an element of a finite level group. This note is the dictionary
that turns such sums into integrals at a finite set $\sT$ of finite places. It shows that the product of the
completions at $\sT$ carries a standard level whose congruence subgroup reads the principal level of the
filling data, that integrality off~$\sT$ and the components at $\sT$ describe the level groups and the unipotent
character, and it proves the \emph{unfolding}: the sum of the amplitudes over a double coset of the unipotent
group, cut off at~$\sT$, is the product of an archimedean Bessel correlation and one coordinate of the local
correlation of Note~T. A second theorem writes a block of amplitude sums with a given link as a sum of limits
of such coset amplitudes as $\sT$ grows. The last section attaches the filling data to the admissible data of
the proof, proves that Bessel systems of the required type exist at the infinite places, and verifies the
standing hypotheses of the filling theorem.

\statement{leancol}{\fontsize{9.6}{13.2}\selectfont\textbf{What this note is.} The artifact \ArtifactName{} is a
body of Lean code: the compact edition of a formal proof, due to L.~Alpöge with the AI system Claude, of the
Riemann hypothesis and of its generalisation to Dirichlet $L$-functions. That headline is a claim. The authors
report that every file compiles and that Lean's kernel accepted the proofs; an independent replay is recorded in
the handout's record of checks when it reports. No mathematician has read the argument through, and nothing here
says that the Riemann hypothesis is proved. This note (Note~G of its unit) treats \UnitInNote{} of the
\UnitMembers{} declarations of the unit \emph{Translates of Poincaré series fill the rows}: the modules on the
passage from global data to the places of~$\sT$, on the filling data of the admissible data and on their
admissibility. Proved here from the inputs: Theorems A to~C and the statements of
§§\ref{s:dict}--\ref{s:unfold}. Theorem~D and Proposition~\ref{p:bessel} are assemblies; the note says which
inputs they assemble and does not reproduce their arithmetic. Taken for granted and stated as Inputs 1 to~6 in
§\ref{s:inputs}: the set-up of the unit, the local theory of Note~T, the standard levels, facts on adeles from
the background library, and the results on the data from Chapter~III and from the units on the columns, the
rows and the discrete series. Every statement was compared with its Lean text, argument by argument. The Lean
text read is \PinText; Appendix~A is the concordance of \ConcDecls{} declarations.}

\sect{The statements}{s:intro}
The objects are recalled in §\ref{s:setting}. In brief: $S$ is filling data over a CM field~$E$,
\m{\GS{S}{E}} the rational unitary group, \m{\GFin{S}} its finite adelic points in the lattice basis, with the
level groups \m{\KgS{S}\le\KS{S}} and the reduction \m{\redS{S}} to the finite group \m{\GRes{S}};
\m{\pos{S}{c}{c'}{\gamma}} is the position of a rational element for two classes, \m{\loc{S}{\sT}} takes the components
at~$\sT$, and \m{\cosamp{S}{\sT}{c}{s}{c'}{\gamma}} is the coset amplitude of the side~$\sigma$ (Definition~\ref{d:cosamp}).

\statement{arith}{\phantomsection\label{t:A}\textbf{Theorem A} (unfolding).\enspace
\emph{In Setting~\ref{s:T}, suppose that \m{\gid{S}^{-1}} is integral off~$\sT$, that $E$ carries Weil data and
that $S$ is admissible. Let $\sigma$ be a side, \m{\gamma_1\in\Hl{\varrho}{\Trat{S}{E}}}, $p$ a sequence of units
of~$E$, and $c$, $c'$ classes such that \m{\intoff{S}{\sT}{\pos{S}{c}{c'}{\weyl{\varrho}{p}\gamma_1}}}. Let
\m{Z\in\IntG{\cL}}, \m{k\in\KS{S}}, and let \m{a,b\in\Cos{Z}} be cosets with
\m{\loc{S}{\sT}(k)^{-1}a_aa_b^{-1}\in\cgp{\cL}}. Put \m{\gamma=\weyl{\varrho}{p}\gamma_1}, let $p_{\sT}$ be the sequence
of the components of $p$ at~$\sT$, and \m{W^{\sigma}=\besfn{\sigma}{S}{\cdot}{\cdot}}. Then}
\begin{align*}&\corr{Z}{p_{\sT}}{\loc{S}{\sT}(\pos{S}{c}{c'}{\gamma_1})^{-1}}_{(a,b)}\cdot\bcorr{dy}{W^{\sigma}}{\rinf{S}(\gamma)}\\
&\qquad=\hdens{\cL}{\varrho}{p_{\sT}}{1}\cdot\gM{\cL}\bigl(\cgp{\cL}\bigr)\cdot\cosamp{S}{\sT}{c}{\redS{S}(k)}{c'}{\gamma}.\end{align*}}
\leanline{\lean{correlationsmuleqcosetAmplitude}}

\statement{arith}{\phantomsection\label{t:B}\textbf{Theorem B} (blocks as limits of coset amplitudes).\enspace
\emph{In Setting~\ref{s:S}, suppose that $E$ carries Weil data and that $S$ is admissible. Let $\sigma$ be a
side, $c$, $c'$ classes, \m{s\in\GRes{S}} and $\Lambda$ a link, such that
\m{\ampsum{\sigma}{S}{\tlink{S}{\Lambda}}{c}{s}{c'}\ne0} and such that for every \m{\gamma\in\terms{S}{c}{s}{c'}} all leading
minors of \m{\Zf{S}{\gamma}} are non-zero. Then there are a unit $t$ of~$E$, an element \m{\gamma_1\in\GS{S}{E}} and
\m{k\in\KS{S}} such that}
\begin{enumerate}[label=\textup{(\roman*)},leftmargin=2.3em,itemsep=1pt]
\item \m{\gamma_1\in\Hl{\varrho}{\Trat{S}{E}}}, \m{\Zf{S}{\gamma_1}_{\varrho\varrho}=t}, \m{\NS{S}^2<|t_\tau|^2} \emph{for every
infinite place~$\tau$,} \m{\redS{S}(k)=s} \emph{and} \m{\Lambda=\link{\weyl{\varrho}{1}\gamma_1}}; \emph{and}
\item \emph{there is a family \m{(L_p)} of vectors of \m{\ell^2(\Bc{S}\times\Bc{S})}, indexed by the maps
\m{p\colon\{0,\dots,\rh{\fs}-1\}\to E^{\times}}, which is summable with sum
\m{\ampsum{\sigma}{S}{\tlink{S}{\Lambda}}{c}{s}{c'}}, such that for every~$p$}
\[\cosamp{S}{\sT}{c}{s}{c'}{\weyl{\varrho}{\pad{p}{t}}\gamma_1}\ \longrightarrow\ L_p\]
\emph{as $\sT$ grows in the directed set of the finite sets of finite places of~$E$.}
\end{enumerate}}
\leanline{\lean{existshasSumcosetAmplitude}}

\statement{arith}{\phantomsection\label{t:C}\textbf{Theorem C} (filling levels exist).\enspace
\emph{In Setting~\ref{s:S}:}
\begin{enumerate}[label=\textup{(\roman*)},leftmargin=2.3em,itemsep=1pt]
\item \emph{Suppose that $E$ carries Weil data, and let $\sT$ be a finite set of finite places of~$E$ that is
stable under complex conjugation. Then there are an integral place structure on \m{\ET{\sT}} and a level $\cL$
of \m{\Tlat{S}{\ET{\sT}}} with \m{\isFL{S}{\sT}{\cL}}.}
\item \emph{For every finite set $\sT_0$ of units of~$E$ there are arbitrarily large finite sets $\sT$ of finite
places, stable under complex conjugation, such that $t$ and $t^{-1}$ are integral off $\sT$ for all
\m{t\in\sT_0}, \m{\gid{S}^{-1}} is integral off~$\sT$, \m{\latB{S}} and its inverse have entries integral
off~$\sT$, and \m{\NS{S}^{-1}} is integral off~$\sT$.}
\item \emph{Every finite adele $a$ is integral off all sufficiently large finite sets $\sT$ of finite places;
and for every \m{x\in\GFin{S}}, \m{\intoff{S}{\sT}{x}} for all sufficiently large~$\sT$.}
\end{enumerate}}
\leanline{\lean{existsisFillingLevel}\sep \lean{existsplacesintegralOff}\sep \lean{eventuallymemintegralOff}\sep
\lean{eventuallyisIntegralOff}}

\statement{arith}{\phantomsection\label{t:D}\textbf{Theorem D} (the filling data of the admissible data).\enspace
\emph{In Setting~\ref{s:D}, let \m{S=\SD{D}} be the filling data of~$D$ (Definition~\ref{d:SD}), and let $K_f$ be
the group of the integral points of the finite adelic unitary group of the standard form on the index set of
the columns. Then:}
\begin{enumerate}[label=\textup{(\roman*)},leftmargin=2.3em,itemsep=1pt]
\item \m{\Colv{D}{0}=\incH{\SD{D}}^{-1}\bigl(\KgS{\SD{D}}\bigr)};
\item \emph{the index of \m{\incH{\SD{D}}^{-1}(\KS{\SD{D}})\cap K_f} in $K_f$ is} \m{(N(D)^{k(D)}-1)^{f(D)}};
\item \m{\#\Cl{\SD{D}}=\#\redim{\SD{D}}{\mathcal H}\cdot(N(D)^{k(D)}-1)^{f(D)}\cdot\mass{E(D)}{k(D)}}, \emph{where}
\m{\mathcal H=\HT{\Tlat{\SD{D}}{\AEf}}}; \emph{and}
\item \emph{\m{\SD{D}} is admissible.}
\end{enumerate}}
\leanline{\lean{columnLevelzeroeqcomap}\sep \lean{relIndexlatStabilizer}\sep \lean{cardclassesfilling}\sep
\lean{instAdmissible} (an instance without a name)}

Theorems A and~B are Inputs 4(f) and~4(e) of the main note of the unit (Note~F, \emph{Translates of Poincaré
series fill the rows}); Theorem~C and the lemmas of §\ref{s:dict} are its Inputs 4(a) to~(c), the lemmas of
§\ref{s:floor} its Input~4(d), and Theorem~D with Definition~\ref{d:SD} and Proposition~\ref{p:bessel} its
Input~9(a).

\sect{The setting}{s:setting}
\begin{conventions}\phantomsection\label{c:conventions}
The natural numbers are \m{0,1,2,\dots}. $E$ is a CM field with its CM type: a number field that is a CM field,
with complex conjugation \m{x\mapsto\bar x}, together with the choice of one embedding at each infinite
place~$\tau$; \m{x_\tau\in\CC} is the image of \m{x\in E} under it, and \m{\ninf{E}} is the number of infinite
places. A \emph{finite place} $w$ of $E$ is a non-zero prime ideal; \m{\F_w} is its residue field, $E_w$ the
completion, $O_w$ its integers, $\bar w$ the conjugate place. \m{\AEf} is the ring of finite adeles, \m{\hO} its
subring of integral adeles, and $\psi_f$ the standard additive character of \m{\AEf}. A finite set $\sT$ of
finite places is \emph{stable under complex conjugation} if it contains $\bar w$ with every~$w$; for a finite
set $\sT$ of finite places, \m{\ET{\sT}=\prod_{w\in\sT}E_w}, \m{x_{\sT}\in\ET{\sT}} is the family of the components
at $\sT$ of a finite adele or of an element of~$E$, and $\tilde y$ is the extension by zero of
\m{y\in\ET{\sT}} to a finite adele. A finite adele is \emph{integral off~$\sT$} if its components at the places
outside $\sT$ are integral, and an element of $E$ if its image is. The field $E$ is \emph{stably unramified} if
it contains a primitive root of unity of order~$5040$, and $E$ \emph{carries Weil data} if it is stably
unramified and a splitting character in the sense of Chapter~III of the handout is given. A \emph{shape} is a
triple \m{\fs=(\varrho,\kappa,k_0)} of a natural number, a finite set and an element of it; \m{\rh{\fs}=\varrho},
\m{\kk{\fs}=\#\kappa}, \m{\nn{\fs}=\rh{\fs}+1}, \m{\rr{\fs}=2\rh{\fs}+\kk{\fs}+1}. A \emph{semilocal field}, an
\emph{integral place structure} on it, with its ring of integers, character and modulus, a \emph{level} $\cL$,
its unit \m{\gL{\cL}} and congruence subgroup \m{\cgp{\cL}}, the group measure \m{\gM{\cL}} and a \emph{standard
level} are as in Note~T of this unit (\emph{The top floor at a finite place}); on \m{\ET{\sT}} the integers are
\m{O_{\sT}=\prod_{w\in\sT}O_w}, the character is written \m{\psiT{\sT}} and the modulus \m{\modT{\sT}}. A \emph{side}
is one of the two words \m{\full}, \m{\cut}; in this note $\sigma$ always denotes a side. $P$ denotes the
parameters of the proof (Chapter~III). $X^*$ is the conjugate transpose of a matrix, \m{\ell^2(X)} the space
\m{\CC^X} with its inner product, conjugate-linear in the first variable. The letters $T$, $F$ of the formal
text for a set of places and a level are written $\sT$ and~$\cL$.
\end{conventions}

\subhead{Forms and groups.} As in Notes F and~T: \m{\Ix{\fs}} is the index set with \m{e_0,\dots,e_{\varrho-1}},
the elements of~$\kappa$, $z$, and \m{e^*_0,\dots,e^*_{\varrho-1}}; the frame is \m{\mathfrak f_a=e_a}
(\m{a<\varrho}), \m{\mathfrak f_\varrho=z}. For a commutative ring $R$ with involution and an $\fs$-form $T$
over~$R$, \m{\JT{T}} is the split Gram matrix, \m{\GT{T}} its unitary group, \m{h(v,v')=v^*\JT{T}v'},
\m{\ZT{T}{x}} the frame Gram matrix with corner entry \m{h(z,xz)}; \m{\UT{T}}, \m{\HT{T}}, \m{\ST{T}=\UT{T}\HT{T}} are
the Bessel unipotent subgroup, the subgroup fixing all $e_a$, $e^*_a$ and~$z$, and the Bessel subgroup;
\m{\Hl{\varrho}{T}} is the subgroup fixing all $e_a$ and $e^*_a$; \m{\sd{u}} is the sum of the entries of $u$ just
above the diagonal of its frame block; and for a sequence $p$ of units, \m{\weyl{\varrho}{p}} is the Weyl element
of length~$\varrho$, which sends $e_a$ to \m{p_ae^*_a} and $e^*_a$ to \m{\bar p_a^{-1}e_a} for \m{a<\varrho}. For
\m{p\colon\{0,\dots,\varrho-1\}\to R^{\times}} and a unit~$t$, \m{\pad{p}{t}} is the sequence
\m{p_0,\dots,p_{\varrho-1},t,1,1,\dots}. For a unit $\pi$, an element $t$ and \m{x\in R^{\kappa}}, the \emph{tail
vector} \m{\tailv{\pi}{t}{x}} is the vector \m{x-(\pi\bar\pi)^{-1}t\,z}, written in the lattice basis below.

\begin{setting}[filling data]\label{s:S}
$E$ is a CM field with its CM type, $\fs$ is a shape, and $S$ is filling data over $E$ of shape~$\fs$, in the
sense of the set-up of this unit (the note \emph{The filling data}). Its components are: a real number
\m{\bS{S}\ge1000}; a prime number \m{\NS{S}\ge5} and a finite idele \m{\piS{S}\in\hO} with
\m{\piS{S}\,\overline{\piS{S}}=\NS{S}}; two primes \m{\qS{S}{0}}, \m{\qS{S}{1}} different from \m{\NS{S}}, with
\m{\gS{S}=\qS{S}{0}\qS{S}{1}}; for \m{i=0,1} a set \m{\ChS{S}{i}} of \m{\ninf{E}} finite places above
\m{\qS{S}{i}}, the \emph{chosen places}, with residue fields of \m{\qS{S}{i}} elements and such that the
conjugate of a chosen place is never chosen; a generator \m{\vpS{S}} of the different of $E$ with
\m{\overline{\vpS{S}}=-\vpS{S}}; an open compact subgroup \m{\KcolS{S}} of the finite adelic unitary group of the
standard form on~$\kappa$, with trivial class automorphism groups; and for every infinite place $\tau$ a
Bessel system \m{\BsysS{S}{\tau}}. Among its inequalities: \m{10^5\le\rh{\fs}},
\m{\nn{\fs}\,[E:\Q]\le2^{2\rh{\fs}}} and \m{1000\,(2\rh{\fs}+2)\le\kk{\fs}}.
\end{setting}

\subhead{The groups of the set-up.} In Setting~\ref{s:S}, for a commutative ring $R$ with involution over the
integers of~$E$: \m{\Trat{S}{R}} is the $\fs$-form that is the identity on $\kappa$ and \m{-\NS{S}} at~$z$, and
\m{\GS{S}{R}=\GT{\Trat{S}{R}}}; \m{\Tlat{S}{R}} is the Gram matrix of the same form in the \emph{lattice basis}, in
which the basis vector of $k_0$ is replaced by a combination of it and of~$z$ with coefficients
\m{\piS{S}^{-1}} and \m{\NS{S}^{-1}}; \m{\latB{S}} is the matrix of this change of basis over \m{\AEf}, and
\m{\GFin{S}=\GT{\Tlat{S}{\AEf}}}. The homomorphism \m{\rfin{S}\colon\GS{S}{E}\to\GFin{S}} is the diagonal embedding
followed by the change of basis; \m{\rinf{S}\colon\GS{S}{E}\to\prod_\tau G_1(\fs)} is the family of the images at
the infinite places, rescaled to the real form \m{G_1(\fs)} of the library by positive numbers, of which
\m{\fscale{S}{\tau}{a}>0} (\m{a\le\varrho}) are the \emph{frame scales}; \m{\incH{S}} embeds the finite adelic points
of the unitary group of the standard form on $\kappa$ into \m{\GFin{S}} as \m{\HT{\Tlat{S}{\AEf}}}. For
\m{\gamma\in\GS{S}{E}}, \m{\Zf{S}{\gamma}=\ZT{\Trat{S}{E}}{\gamma}}, \m{\tdisc{S}{\gamma}} is the \emph{term discriminant},
equal to \m{d\bar d\,(t\bar t-\NS{S}^2)} when \m{\Zf{S}{\gamma}} is diagonal with entries \m{p_0,\dots,p_{\varrho-1},t}
and \m{d=p_0\cdots p_{\varrho-1}}; \m{\link{\gamma}} is the \emph{link} of~$\gamma$, a subspace attached to $\gamma$ in
the set-up, and for a link~$\Lambda$,
\m{\tlink{S}{\Lambda}=\{\gamma\in\GS{S}{E}:\ \tdisc{S}{\gamma}\ne0,\ \link{\gamma}=\Lambda\}}.

For a finite idele~$u$, \m{\KA{S}{u}} is the group of the \m{x\in\GFin{S}} that have entries in \m{\hO}, as has
$x^{-1}$, and for which \m{u^{-1}(x-1)} has entries in~\m{\hO}. Then \m{\KS{S}=\KA{S}{1}}, and with \m{\gid{S}}, the
integer \m{\gS{S}} as a finite idele, \m{\KgS{S}=\KA{S}{\gid{S}}} is the kernel of the reduction
\m{\redS{S}\colon\KS{S}\to\GRes{S}} modulo the chosen places, where
\m{\GRes{S}=\prod_{i}\prod_{w\in\ChS{S}{i}}\GL_{\Ix{\fs}}(\F_w)}; for a subgroup $B$ of \m{\GFin{S}},
\m{\redim{S}{B}=\redS{S}(B\cap\KS{S})}. The character \m{\psiS{S}(x)=\psi_f(-x/\vpS{S})} of \m{\AEf} is self-dual:
\m{\psiS{S}(ra)=1} for all \m{r\in\hO} exactly if \m{a\in\hO}; and
\m{\uchar{S}{u}=\psiS{S}\bigl(\gid{S}^{-1}\sd{u}\bigr)} for \m{u\in\GFin{S}}. For a finite set $\sT$ of finite places,
\m{\charT{S}{\sT}(y)=\psiS{S}(\tilde y)} is the character of \m{\ET{\sT}} of the set-up. \m{\Cl{S}} is the finite set
of the double cosets of the finite adelic unitary group of $\kappa$ by its rational points and \m{\KcolS{S}};
each class $c$ has a chosen representative $y_c$, and \m{\hc{S}{c}=\incH{S}(y_c)}. For classes $c$, $c'$ and
\m{\gamma\in\GS{S}{E}},
\[\pos{S}{c}{c'}{\gamma}=\hc{S}{c}^{-1}\,\rfin{S}(\gamma)\,\hc{S}{c'}\in\GFin{S},\]
and for \m{s\in\GRes{S}} the set of \emph{terms over} \m{(c,s,c')} is
$\terms{S}{c}{s}{c'}=\{\gamma\in\GS{S}{E}:\ \pos{S}{c}{c'}{\gamma}\in\KS{S},\ \redS{S}(\pos{S}{c}{c'}{\gamma})=s\}$.
\m{\clU{S}{c}} is the group of the rational \m{u\in\UT{\Trat{S}{E}}} with \m{\pos{S}{c}{c}{u}\in\KgS{S}}.

\subhead{Amplitudes.} \m{\Bc{S}} is the set of the families \m{(i_\tau)_\tau} of indices of the Bessel systems. For
a side $\sigma$, a point $y$ of \m{\prod_\tau G_1(\fs)} and \m{i\in\Bc{S}}, \m{\besfn{\sigma}{S}{y}{i}} is the product
over $\tau$ of the $i_\tau$-th coordinates at $y_\tau$ of the full or the cut Bessel function of
\m{\BsysS{S}{\tau}}; \m{\bker{\sigma}{S}{\gamma}{y}\in\ell^2(\Bc{S}\times\Bc{S})}, the \emph{Bessel kernel}, has the
coordinates \m{\besfn{\sigma}{S}{y}{i}\,\overline{\besfn{\sigma}{S}{\rinf{S}(\gamma)^{-1}y}{i'}}}. With the cutoffs
\m{\cutoff{S}{c}\ge0} of the set-up, functions on \m{\prod_\tau G_1(\fs)}, the \emph{amplitude} of
\m{\gamma\in\GS{S}{E}} over $c$, $c'$ and the \emph{amplitude sum} of a set \m{\Xi\subseteq\GS{S}{E}} are
\begin{align*}\amp{\sigma}{S}{c}{c'}{\gamma}&=\int\cutoff{S}{c}(y)\,\cutoff{S}{c'}\bigl(\rinf{S}(\gamma)^{-1}y\bigr)\,\bker{\sigma}{S}{\gamma}{y}\,dy,\\
\ampsum{\sigma}{S}{\Xi}{c}{s}{c'}&=\sum_{\gamma\in\Xi\cap\terms{S}{c}{s}{c'}}\amp{\sigma}{S}{c}{c'}{\gamma},\end{align*}
with $dy$ the Haar measure of the set-up; \m{\ampabs{\sigma}{S}{c}{c'}{\gamma}} is the same integral with the norm of
the Bessel kernel. The \emph{order} is \m{\uord{S}=\gS{S}^{\ninf{E}\rh{\fs}(2\kk{\fs}+2\rh{\fs}+1)}}. $S$ is
\emph{admissible} if the standing hypotheses of the set-up hold: \m{\KcolS{S}=\incH{S}^{-1}(\KgS{S})}, the
reduction \m{\redS{S}} is surjective, and the inequalities listed in the proof of Theorem~D.

\subhead{The local objects of Note~T.} For a finite set $\sT$ of finite places and a unit $u$ of \m{\ET{\sT}}, \m{\KT{S}{\sT}{u}} is the group of the
\m{y\in\GT{\Tlat{S}{\ET{\sT}}}} that have entries in $O_{\sT}$, as has $y^{-1}$, and for which \m{u^{-1}(y-1)} has
entries in~$O_{\sT}$. For a level $\cL$ of \m{\Tlat{S}{\ET{\sT}}}: \m{\cgp{\cL}=\KT{S}{\sT}{\gL{\cL}}};
\m{\phs{g}{y}} is the finite Bessel function; \m{\IntG{\cL}} the set of the intermediate groups $Z$, with
\m{\cgp{\cL}\le Z\le\KT{S}{\sT}{1}}; \m{\Cos{Z}} the set of the cosets of \m{\cgp{\cL}} in $Z$, each coset $c$ with a
representative $a_c$; \m{\bvec{Z}{y}} the finite Bessel vector, with the coordinate
\m{\phs{\gL{\cL}}{y\,a_c^{-1}}} at~$c$; \m{\umass{\cL}=\prod_{l<\varrho}\fmod{\cL}{l}{\gL{\cL}}} the unipotent mass, where
\m{\fmod{\cL}{l}{g}=\modT{\sT}(g)^{\kk{\fs}+1+2(\varrho-l-1)+1/2}}; \m{\hdens{\cL}{\varrho}{p}{1}} the head density; and
for a sequence $p$ of units of \m{\ET{\sT}} and \m{l\in\GT{\Tlat{S}{\ET{\sT}}}}, the correlation is
\m{\corr{Z}{p}{l}=\hdens{\cL}{\varrho}{p}{1}\cdot\bcorr{\gM{\cL}}{\bvec{Z}{\cdot}}{\weyl{\varrho}{p}\,l^{-1}}}, where for a
measure $\mu$, a vector-valued function $W$ and a group element $\gamma$ the \emph{Bessel correlation}
\m{\bcorr{\mu}{W}{\gamma}} is the tensor with the coordinates \m{\int W(y)_i\,\overline{W(\gamma^{-1}y)_{i'}}\,d\mu(y)}.
\m{\etab{j}} is the convolution of the first $j$ Heisenberg measures.

\begin{definition}[integrality off $\sT$ and the components at $\sT$]\label{d:off}
In Setting~\ref{s:S}, let $\sT$ be a finite set of finite places of~$E$.
\begin{enumerate}[label=\textup{(\alph*)},leftmargin=2.3em,itemsep=1pt]
\item For \m{x\in\GFin{S}} we write \m{\intoff{S}{\sT}{x}}, and say that $x$ is \emph{integral off~$\sT$}, if the
entries of $x$ and of $x^{-1}$ are integral off~$\sT$.
\item If $\sT$ is stable under complex conjugation, \m{\loc{S}{\sT}\colon\GFin{S}\to\GT{\Tlat{S}{\ET{\sT}}}} is the
homomorphism that takes the components at $\sT$ of every entry.
\end{enumerate}
\end{definition}
\leanline{\lean{IsIntegralOff}\sep \lean{toPlaces}}

\begin{definition}[filling levels]\label{d:FL}
In Setting~\ref{s:S}, let $\sT$ be a finite set of finite places of $E$ that is stable under complex
conjugation, suppose that \m{\ET{\sT}} carries an integral place structure, and let
\[\cL\in\Levels{\{1,\dots,\kk{\fs}-6\rh{\fs}-13\}}{\Tlat{S}{\ET{\sT}}}{3\rh{\fs}+6}\]
be a level, with \m{3\rh{\fs}+6} further hyperbolic pairs and an orthonormal part of \m{\kk{\fs}-6\rh{\fs}-13}
vectors. Then \m{\isFL{S}{\sT}{\cL}} (\emph{$\cL$ is a filling level for $S$ at~$\sT$}) means: $\cL$ is a standard
level, \m{\IsStd{\cL}}; the character of the place structure is \m{\psiT{\sT}=\charT{S}{\sT}}; and
\m{\gL{\cL}=\gid{S}_{\sT}}.
\end{definition}
\leanline{\lean{IsFillingLevel}}

\begin{setting}[a filling level]\label{s:T}
In Setting~\ref{s:S}, with $E$, $\fs$ and $S$ as there: $\sT$ is a finite set of finite places of $E$ that is
stable under complex conjugation; \m{\ET{\sT}} carries an integral place structure; and $\cL$ is a level of
\m{\Tlat{S}{\ET{\sT}}} with \m{\isFL{S}{\sT}{\cL}}, a filling level.
\end{setting}

\begin{definition}[coset amplitudes]\label{d:cosamp}
In Setting~\ref{s:S}, let $\sT$ be a finite set of finite places of~$E$, $\sigma$ a side, $c$, $c'$ classes,
\m{s\in\GRes{S}} and \m{\gamma\in\GS{S}{E}}. The \emph{coset amplitude} is
\[\cosamp{S}{\sT}{c}{s}{c'}{\gamma}=\sum_{(u,u')}\amp{\sigma}{S}{c}{c'}{u\gamma u'^{-1}}\ \in\ \ell^2(\Bc{S}\times\Bc{S}),\]
the sum over the pairs \m{(u,u')} of elements of \m{\UT{\Trat{S}{E}}} with \m{\intoff{S}{\sT}{\pos{S}{c}{c}{u}}} and
\m{u\gamma u'^{-1}\in\terms{S}{c}{s}{c'}}. The side is an instance argument of the formal definition; it is the
superscript $\sigma$ of the notation, and every statement of this note says which side it is.
\end{definition}
\leanline{\lean{cosetAmplitude}}

\begin{setting}[the admissible data]\label{s:D}
\m{P=(N_0,d)} are the parameters of the proof and $D$ is admissible data for~$P$ (Chapter~III of the handout).
Attached to $D$ are the CM field \m{E(D)} with its CM type, which carries Weil data; the ranks \m{m(D)},
\m{r(D)}, \m{k(D)} with \m{r(D)+1=k(D)+m(D)} and \m{m(D)=2\varrho(D)+2}; the number \m{f(D)} of the infinite places
of \m{E(D)}; the third scale \m{b(D)}; the level \m{N(D)}, a prime, and a finite idele \m{\hat\pi(D)} with
\m{\hat\pi(D)\overline{\hat\pi(D)}=N(D)}; the primes \m{g_0(D)}, \m{g_1(D)} with their chosen places; the generator
\m{\varpi(P)} of the different; from the units on the columns and on the seesaw, the index set \m{X_0(D)} of
the columns, with \m{k(D)} elements (Lemma~\ref{l:small}), the level \m{\Colv{D}{0}} of the columns, the shape
\m{\fs(D)=(\varrho(D),X_0(D),k_0)} with $k_0$ the last index, the finite group \m{\sGD{D}}, and for every infinite
place $\tau$ the doubled row parameter \m{2\varkappa_D(\tau,\cdot)} and the column weight
\m{\mathsf{w}^{\mathrm{col}}_D(\tau)}, a weight of the unitary group of the columns.
\end{setting}

\sect{What is taken for granted}{s:inputs}
The following statements are used and not proved here. Inputs 1 to~3 belong to the set-up of this unit (the
notes \emph{The filling data}, \emph{Box partitions of arithmetic unipotent groups} and \emph{The energy of
amplitude sums}) and, for the big cell, to the background library; Input~4 to Note~T and to the note
\emph{Poincaré series of cut Bessel functions} of this unit; Input~5 to the background library; Input~6 to
Chapter~III and to the units on the columns, the rows, the real place and the discrete series. They are stated
from the dossier's readings of their quoted Lean statements. Where a classical theorem stands behind an input,
its reference is to be confirmed, except where another note is cited for it.

\statement{est}{\phantomsection\label{i:1}\textbf{Input 1} (levels and the unipotent character; set-up of this
unit).\enspace In Setting~\ref{s:S}, with \m{\mathcal U=\UT{\Tlat{S}{\AEf}}} and \m{\mathcal H=\HT{\Tlat{S}{\AEf}}}:
\begin{enumerate}[label=\textup{(\alph*)},leftmargin=2.3em,itemsep=1pt]
\item The kernel of \m{\redS{S}} is \m{\KA{S}{\gid{S}}}, as said in §\ref{s:setting}; \m{\gid{S}\in\hO}; the split Gram
matrix of \m{\Tlat{S}{\AEf}} and its inverse have entries in~\m{\hO}; and the components of \m{\Tlat{S}{\AEf}} at $\sT$
form \m{\Tlat{S}{\ET{\sT}}}, the lattice form for the integer \m{\NS{S}} of~\m{\ET{\sT}}.
\item \m{\psiS{S}(ra)=1} for all \m{r\in\hO} exactly if \m{a\in\hO}; and \m{\charT{S}{\sT}(x_{\sT})=\psiS{S}(x)} for every
finite adele $x$ that is integral off~$\sT$.
\item For a rational \m{u\in\UT{\Trat{S}{E}}} and a class~$c$: \m{\rfin{S}(u)\in\mathcal U} and
\m{\pos{S}{c}{c}{u}\in\mathcal U}. If \m{u,v\in\mathcal U} and $h$ is in the image of \m{\incH{S}} with
\m{u^{-1}h^{-1}vh\in\KgS{S}}, then \m{\uchar{S}{v}=\uchar{S}{u}}.
\item \emph{(Rational approximation.)} For a class $c$ and \m{u\in\mathcal U} there is a rational
\m{u_0\in\UT{\Trat{S}{E}}} with \m{u^{-1}\pos{S}{c}{c}{u_0}\in\KgS{S}} and \m{\uchar{S}{\rfin{S}(u_0)}=\uchar{S}{u}}.
\item For a form $T$ as in Note~T over \m{\AEf} or over \m{\ET{\sT}}, a unit~$g$, \m{u\in\UT{T}} and \m{h\in\HT{T}}: if
$uh$ lies in the congruence subgroup modulo~$g$, so does~$h$. Every element of \m{\ST{T}} is such a product~$uh$.
\end{enumerate}}

\statement{est}{\phantomsection\label{i:2}\textbf{Input 2} (amplitudes; set-up of this unit).\enspace In
Setting~\ref{s:S}, let $\sigma$ be a side and $c$, $c'$ classes.
\begin{enumerate}[label=\textup{(\alph*)},leftmargin=2.3em,itemsep=1pt]
\item \emph{(Partition of unity.)} If $E$ is stably unramified and $S$ admissible, then
\m{\sum_{u\in\clU{S}{c}}\cutoff{S}{c}\bigl(\rinf{S}(u)\,y\bigr)=1} for every point~$y$.
\item \emph{(The kernel under unipotent elements.)} For rational \m{u,u'\in\UT{\Trat{S}{E}}}, \m{\gamma\in\GS{S}{E}} and
a point~$y$: \m{\bker{\sigma}{S}{u\gamma u'^{-1}}{\rinf{S}(u)\,y}=\uchar{S}{\rfin{S}(u)}^{-1}\,\uchar{S}{\rfin{S}(u')}\cdot
\bker{\sigma}{S}{\gamma}{y}}.
\item \emph{(Summability.)} If $E$ is stably unramified and $S$ admissible, then for every
\m{\Xi\subseteq\GS{S}{E}} and \m{s\in\GRes{S}} the family of the amplitudes \m{\amp{\sigma}{S}{c}{c'}{\gamma}},
\m{\gamma\in\Xi\cap\terms{S}{c}{s}{c'}}, is summable, and \m{\sum_{\gamma\in\terms{S}{c}{s}{c'}}\ampabs{\sigma}{S}{c}{c'}{\gamma}<\infty}.
\end{enumerate}}

\statement{est}{\phantomsection\label{i:3}\textbf{Input 3} (the big cell and links; background library and
set-up).\enspace Let $T$ be an $\fs$-form over a field with involution whose split Gram matrix is hermitian
and invertible; here \m{T=\Trat{S}{E}}.
\begin{enumerate}[label=\textup{(\alph*)},leftmargin=2.3em,itemsep=1pt]
\item If all leading minors of \m{\ZT{T}{\gamma}} are non-zero, there are \m{u,u'\in\UT{T}}, a sequence $p$ of units
and \m{\gamma_1\in\Hl{\varrho}{T}} with \m{\gamma=u^{-1}\,\weyl{\varrho}{p}\gamma_1\,u'}.
\item If \m{v_1\,\weyl{\varrho}{p}\gamma_1\,v_2^{-1}=v_1'\,\weyl{\varrho}{q}\gamma_1\,v_2'^{-1}} with
\m{v_1,v_2,v_1',v_2'\in\UT{T}} and \m{\gamma_1\in\Hl{\varrho}{T}}, then \m{p_a=q_a} for \m{a<\varrho}, \m{v_1=v_1'} and
\m{v_2=v_2'}.
\item For \m{\gamma_1\in\Hl{\varrho}{T}}, \m{\ZT{T}{\weyl{\varrho}{p}\gamma_1}} is the diagonal matrix with the entries
\m{p_0,\dots,p_{\varrho-1}} and \m{h(z,\gamma_1z)}. For \m{u,u'\in\UT{T}}, the leading minors of \m{\ZT{T}{u^{-1}\gamma u'}}
are those of \m{\ZT{T}{\gamma}}, the term discriminants of \m{u^{-1}\gamma u'} and $\gamma$ agree, and
\m{\link{u^{-1}\gamma u'}=\link{\gamma}}. For \m{\gamma_1,\gamma_2\in\Hl{\varrho}{T}} and sequences $p$, $q$ of units:
\m{\link{\weyl{\varrho}{p}\gamma_1}=\link{\weyl{\varrho}{q}\gamma_2}} exactly if \m{\gamma_1=\gamma_2}.
\end{enumerate}}

\statement{est}{\phantomsection\label{i:4}\textbf{Input 4} (the local theory; Note~T and the note on standard
levels).\enspace Let $A$ be a semilocal field with an involution and an integral place structure, $T$ a form
and $\cL$ a level as in Note~T, with integers $O_A$, character $\psi_A$ and modulus \m{\operatorname{mod}_A}.
\begin{enumerate}[label=\textup{(\alph*)},leftmargin=2.3em,itemsep=1pt]
\item \emph{(Note~T, Lemma on the finite Bessel function.)} For a unit $g$ and \m{y\in\GT{T}}: if \m{\phs{g}{y}\ne0}
there is \m{u\in\UT{T}} with $u^{-1}y$ in the congruence subgroup modulo~$g$; for every such~$u$,
\m{\phs{g}{y}=\psi_A(g^{-1}\sd{u})}; \m{\phs{g}{yk}=\phs{g}{y}} for $k$ in that subgroup; and
\m{\phs{g}{uy}=\psi_A(g^{-1}\sd{u})\phs{g}{y}} for \m{u\in\UT{T}}. Moreover \m{\sd{h}=0} for \m{h\in\HT{T}}, and
\m{\sd{uu'}=\sd{u}+\sd{u'}} for \m{u,u'\in\ST{T}}.
\item \emph{(Standard levels.)} If \m{\IsStd{\cL}}: \m{\gM{\cL}} is a Haar measure, invariant under left and right
translations and finite on the group of the integral isometries; \m{\etab{\varrho}} is carried by \m{\UT{T}} and
invariant under right translation by its elements; and for a unit \m{g\in O_A} with \m{\bar g=g}, the congruence
subgroup modulo $g$ has \m{\etab{\varrho}}-measure \m{\prod_{l<\varrho}\fmod{\cL}{l}{g}}.
\item \emph{(Lattice levels.)} Let \m{N\in O_A} with \m{\bar N=N} be such that the split Gram matrix of the
lattice form of the shape $\fs$ for $N$ is a hermitian unit, $g$ a unit of $A$ with \m{g\in O_A}, \m{\bar g=g}, and
$\iota$ a finite set, $m'$ a natural number with \m{\kk{\fs}=\#\iota+2m'+1}. Then that form has a level
in \m{\Levels{\iota}{\cdot}{m'}} with unit $g$ that is a \emph{lattice level}. A level in
\m{\Levels{\iota}{\cdot}{m+1}} is standard exactly if it is a lattice level, \m{\zeta_A(m)^{\nn{\fs}}\le1+2^{-\nn{\fs}}}
for the ideal zeta function of~$O_A$, \m{m<9\nn{\fs}}, \m{13\nn{\fs}\le\#\iota}, and the unitary group has a Haar
measure that is right invariant.
\end{enumerate}}

\statement{est}{\phantomsection\label{i:5}\textbf{Input 5} (adeles; background library).\enspace Let $\sT$ be a
finite set of finite places of~$E$.
\begin{enumerate}[label=\textup{(\alph*)},leftmargin=2.3em,itemsep=1pt]
\item \emph{(Extension by zero.)} The map \m{y\mapsto\tilde y} from \m{\ET{\sT}} to \m{\AEf} is additive and
multiplicative, \m{(\tilde y)_{\sT}=y}, $\tilde y$ is integral off~$\sT$, and \m{\tilde{\bar y}=\bar{\tilde y}} if $\sT$ is
stable under complex conjugation. A finite adele lies in \m{\hO} exactly if it is integral off $\sT$ and its
components at $\sT$ are integral.
\item \emph{(Indices and archimedean absolute values.)} With the generalised index \m{\gix{\cdot}} of
\m{\ET{\sT}}: for \m{x\in E^{\times}} integral off~$\sT$, \m{\gix{x_{\sT}}^{1/2}\le\prod_\tau|x_\tau|}.
\item \emph{(Zeta functions.)} The ideal zeta function of $O_{\sT}$ is at most that of~\m{\hO}; and if
\m{n\,[E:\Q]\le2^{v}}, then \m{\zeta_{\hO}(u+v+4)^{n}\le1+2^{-u}}.
\item \emph{(Unimodular lattices.)} A hermitian matrix $b$ over \m{\AEf} that has, with its inverse, entries in
\m{\hO} is \m{k^*bk=1} for some matrix $k$ with $k$, $k^{-1}$ integral. The image of a Haar measure of a locally
compact group, restricted to an open subgroup, under a continuous homomorphism that is surjective on that
subgroup is a Haar measure of the image; and \m{\ET{\sT}} has a unit $\delta$ with \m{\bar\delta=-\delta},
\m{\delta,\delta^{-1}\in O_{\sT}}, and an element \m{\tau_0\in O_{\sT}} with \m{\tau_0+\bar\tau_0=1}.
\item \emph{(Finite groups.)} The stabiliser of a non-zero vector of $k^n$ in \m{\GL_n(k)}, for a finite field
$k$ and \m{n\ge1}, has index \m{\#k^{n}-1}.
\end{enumerate}}

\statement{est}{\phantomsection\label{i:6}\textbf{Input 6} (the data; Chapter~III and other units).\enspace In
Setting~\ref{s:D}:
\begin{enumerate}[label=\textup{(\alph*)},leftmargin=2.3em,itemsep=1pt]
\item \emph{(Chapter~III.)} The logarithmic bounds of the admissible data, in particular
\m{b(D)^{19}\le g_0(D)g_1(D)\le b(D)^{23}}, the window for \m{\log N(D)}, \m{k(D)\le a(D)}, \m{\pi\le b(D)}, and the
inequalities of Setting~\ref{s:S} for \m{\fs(D)} and \m{b(D)}; the primes \m{g_i(D)} are different from \m{N(D)}, an
integral adele is divisible by \m{g_0(D)g_1(D)} exactly if it and its conjugate vanish modulo all chosen
places, the conjugate of a chosen place is not chosen, and there are \m{f(D)} chosen places above each
\m{g_i(D)}; \m{\varpi(P)} generates the different and \m{\overline{\varpi(P)}=-\varpi(P)}.
\item \emph{(Columns.)} The class automorphism groups of \m{\Colv{D}{0}} are trivial, so that the number of its
classes is a class sum, which for a subgroup of $K_f$ is its index times the class sum of~$K_f$; the class sum
of $K_f$ is \m{\mass{E(D)}{k(D)}}; \m{\Colv{D}{0}} is the group of the elements of $K_f$ that reduce to $1$ at all
chosen places and whose row at the last index is congruent to that of the identity modulo every place
dividing~\m{\hat\pi(D)};
there are \m{f(D)} such places, each with \m{N(D)} elements in its residue field, none conjugate to another.
\item \emph{(The filling data fit.)} The filling data \m{\SD{D}} of Definition~\ref{d:SD} fit the space
\m{Y_0(D)} of the rows in the sense of the note on Poincaré series, and every family of residue matrices at the
chosen places is the reduction of an integral element of the finite adelic unitary group; consequently
\m{\redS{\SD{D}}} is surjective. Under \m{\incH{S}}, an element $y$ lies in \m{\incH{S}^{-1}(\KS{S})} exactly if it is
integral and its row at $k_0$ is congruent to that of the identity modulo every place dividing~\m{\piS{S}}, and
the reduction of \m{\incH{S}(y)} at a chosen place is $1$ exactly if that of $y$ is (set-up of the unit).
\item \emph{(Bessel functions; units on the discrete series, on the real place and on the exchange.)} For a
Harish-Chandra parameter with a large gap: the formal degree is positive and is given by Weyl's dimension
polynomial (the reference and its relation to the formal statement are given in Note~F, Input~8(g)); the kernel
of the discrete series in a frame is idempotent; there is a Bessel function of the type of an irreducible
unitary representation of the unitary group of the columns whose matrix coefficient is the coefficient of the
lowest type, with the bounds of the Laplace method for steep parameters; a Bessel function with these bounds
has a cut pair; the full function of a cut pair satisfies Schur bounds; and a cut Bessel function has band
limits.
\item \emph{(Sizes of the counts.)} A bound for \m{\#\Bc{\SD{D}}}, the logarithm of the degree \m{\bdeg{\SD{D}}}
within an explicit interval, and \m{|\log(\mass{E(D)}{i+1}/\mass{E(D)}{i})-\cdots|} bounded for
\m{k(D)\le i<r(D)}, as the formal text states them (the note on the degree bounds of this unit and the library
on masses).
\item \emph{(Line coordinates; unit on the rows.)} With the notation of the note \emph{The height of a ratio of
lift volumes}: the polynomial of the line coordinates \m{\linePolyCoords{j}{i}{Y}} in the weight of the first
columns is the product $\Delta$ of the stacked determinants of the matrices \m{\detRows{j}{i}{Y}{\tau}} along
\m{\posEquiv{j}}, read as a polynomial in the first columns with coefficients polynomials in the other columns.
\end{enumerate}}

\sect{The dictionary between adelic and local objects}{s:dict}
\begin{lemma}[lifting from the places of $\sT$]\label{l:lift}
In Setting~\ref{s:S}, let $\sT$ be a finite set of finite places of~$E$ that is stable under complex
conjugation. For every \m{y\in\GT{\Tlat{S}{\ET{\sT}}}} there is \m{x\in\GFin{S}} with \m{\loc{S}{\sT}(x)=y} and
\m{\intoff{S}{\sT}{x}}, such that \m{x\in\UT{\Tlat{S}{\AEf}}} if \m{y\in\UT{\Tlat{S}{\ET{\sT}}}}, and
\m{x\in\HT{\Tlat{S}{\AEf}}} if \m{y\in\HT{\Tlat{S}{\ET{\sT}}}}.
\end{lemma}
\begin{proof}
Let \m{e=y-1} and \m{x=1+\tilde e}, the extension by zero taken entry by entry. Let $J$ be the split Gram matrix
of \m{\Tlat{S}{\AEf}} and $J_{\sT}$ that of \m{\Tlat{S}{\ET{\sT}}}, the matrix of the components of~$J$. Extension by
zero is multiplicative and commutes with conjugation (Input~\ref{i:5}(a)), and \m{\tilde a\,b=\widetilde{a\,b_{\sT}}}
for a finite adele~$b$. So
\[x^*Jx-J=\tilde e^*J+J\tilde e+\tilde e^*J\tilde e=\widetilde{e^*J_{\sT}+J_{\sT}e+e^*J_{\sT}e}=\widetilde{y^*J_{\sT}y-J_{\sT}}=0,\]
and \m{x\in\GFin{S}}. Its components at $\sT$ are those of~$y$; off $\sT$ it is the identity, and its inverse is
\m{1+\widetilde{y^{-1}-1}}, so \m{\intoff{S}{\sT}{x}}. The groups $\mathcal U$ and $\mathcal H$ are defined by
conditions that certain entries are $0$ or~$1$; these pass from $y$ to~$x$.
\end{proof}
\leanline{\lean{existstoPlaceseq}}

\begin{lemma}[level groups at $\sT$]\label{l:levels}
In Setting~\ref{s:S}, let $\sT$ be a finite set of finite places of $E$ that is stable under complex
conjugation, and \m{x\in\GFin{S}}.
\begin{enumerate}[label=\textup{(\alph*)},leftmargin=2.3em,itemsep=1pt]
\item For a finite idele $u$ such that \m{u^{-1}} is integral off~$\sT$: \m{x\in\KA{S}{u}} exactly if
\m{\intoff{S}{\sT}{x}} and \m{\loc{S}{\sT}(x)\in\KT{S}{\sT}{u_{\sT}}}.
\item \m{x\in\KS{S}} exactly if \m{\intoff{S}{\sT}{x}} and \m{\loc{S}{\sT}(x)\in\KT{S}{\sT}{1}}.
\end{enumerate}
Suppose moreover that we are in Setting~\ref{s:T}, with a filling level~$\cL$, and that \m{\gid{S}^{-1}} is integral
off~$\sT$.
\begin{enumerate}[label=\textup{(\alph*)},leftmargin=2.3em,itemsep=1pt,resume]
\item \m{x\in\KgS{S}} exactly if \m{\intoff{S}{\sT}{x}} and \m{\loc{S}{\sT}(x)\in\cgp{\cL}}.
\item For \m{k,l\in\KS{S}}: \m{\redS{S}(k)=\redS{S}(l)} exactly if \m{\loc{S}{\sT}(k)^{-1}\loc{S}{\sT}(l)\in\cgp{\cL}}.
\end{enumerate}
\end{lemma}
\begin{proof}
(a) A finite adele is integral exactly if it is integral off $\sT$ and its components at $\sT$ are integral
(Input~\ref{i:5}(a)). If the entries of $x$ are integral off~$\sT$, so are those of \m{u^{-1}(x-1)}, because
\m{u^{-1}} is. So the three conditions that define \m{\KA{S}{u}} split into \m{\intoff{S}{\sT}{x}} and the three
conditions at~$\sT$, which define \m{\KT{S}{\sT}{u_{\sT}}}. (b) is (a) for \m{u=1}. (c) is (a) for \m{u=\gid{S}}, as
\m{\KgS{S}=\KA{S}{\gid{S}}} (Input~\ref{i:1}(a)) and \m{\cgp{\cL}=\KT{S}{\sT}{\gL{\cL}}} with \m{\gL{\cL}=\gid{S}_{\sT}}.
(d) The reductions agree exactly if \m{k^{-1}l\in\KgS{S}}; as \m{k^{-1}l\in\KS{S}} is integral off $\sT$ by~(b),
this is the condition of~(c).
\end{proof}
\leanline{\lean{memunitaryLeveliff}\sep \lean{memlatStabilizeriff}\sep \lean{memprincipalLeveliff}\sep
\lean{reductioneqiff}}

\begin{lemma}[the unipotent character at $\sT$]\label{l:char}
In Setting~\ref{s:T}, suppose that \m{\gid{S}^{-1}} is integral off~$\sT$, and let \m{x\in\GFin{S}} with
\m{\intoff{S}{\sT}{x}}. Then
\m{\psiT{\sT}\bigl(\gL{\cL}^{-1}\,\sd{\loc{S}{\sT}(x)}\bigr)=\uchar{S}{x}}.
\end{lemma}
\begin{proof}
The left side is \m{\charT{S}{\sT}(a_{\sT})} for the finite adele \m{a=\gid{S}^{-1}\sd{x}}, because
\m{\psiT{\sT}=\charT{S}{\sT}} and \m{\gL{\cL}=\gid{S}_{\sT}}. As $a$ is integral off~$\sT$, this is
\m{\psiS{S}(a)=\uchar{S}{x}} by Input~\ref{i:1}(b).
\end{proof}
\leanline{\lean{psisimpleRootSumequnipotentChar}}

\begin{lemma}[moduli]\label{l:mod}
In Setting~\ref{s:S}, let $\sT$ be a finite set of finite places of $E$ that is stable under complex
conjugation, and suppose that \m{\ET{\sT}} carries an integral place structure.
\begin{enumerate}[label=\textup{(\alph*)},leftmargin=2.3em,itemsep=1pt]
\item For a unit $x$ of $E$ such that $x$ and $x^{-1}$ are integral off~$\sT$:
$\modT{\sT}(x_{\sT})^{1/2}=\bigl(\prod_\tau|x_\tau|\bigr)^{-1}$.
\end{enumerate}
Suppose moreover that we are in Setting~\ref{s:T}, with a filling level~$\cL$, and that \m{\gid{S}^{-1}} is integral
off~$\sT$.
\begin{enumerate}[label=\textup{(\alph*)},leftmargin=2.3em,itemsep=1pt,resume]
\item \m{\modT{\sT}(\gL{\cL})=\gS{S}^{-2\ninf{E}}}.
\item \m{\uord{S}\cdot\umass{\cL}=1}.
\end{enumerate}
\end{lemma}
\begin{proof}
(a) By Input~\ref{i:5}(b) for $x$ and for $x^{-1}$, \m{\gix{x_{\sT}}^{1/2}\le\prod_\tau|x_\tau|} and
\m{\gix{x_{\sT}}^{-1/2}\le\prod_\tau|x_\tau|^{-1}}; so \m{\gix{x_{\sT}}^{1/2}=\prod_\tau|x_\tau|}, and
\m{\modT{\sT}(x_{\sT})=\gix{x_{\sT}}^{-1}}. (b) is (a) for the integer \m{x=\gS{S}}, whose absolute value at every
infinite place is \m{\gS{S}}; its inverse is integral off $\sT$ by hypothesis. (c) By (b),
\[\umass{\cL}=\prod_{l<\varrho}\gS{S}^{-2\ninf{E}\,(\kk{\fs}+1+2(\varrho-l-1)+1/2)}=\gS{S}^{-\ninf{E}\,\varrho\,(2\kk{\fs}+2\varrho+1)},\]
since \m{\sum_{l<\varrho}(\kk{\fs}+\tfrac32+2(\varrho-l-1))=\varrho(\kk{\fs}+\varrho+\tfrac12)}; this is \m{\uord{S}^{-1}}.
\end{proof}
\leanline{\lean{mdevalPlacesrpowhalf}\sep \lean{mdg}\sep \lean{gpowmulunipotentMass}}

\begin{proof}[Proof of Theorem C]
\emph{(iii)} A finite adele is integral at all but finitely many places; so are the finitely many entries of
$x$ and of~$x^{-1}$.

\emph{(ii)} Given a finite set $\sT'$ of finite places, the finitely many finite adeles in question are
integral off every sufficiently large set, by~(iii); choose such a set that contains $\sT'$, and add the
conjugates of its places.

\emph{(i)} Write \m{A=\ET{\sT}}. \emph{The place structure.} By Input~\ref{i:5}(d) there are a unit $\delta$ with
\m{\bar\delta=-\delta} and an element $\tau_0$ with \m{\tau_0+\bar\tau_0=1}, integral with integral inverse of~$\delta$.
The data of the structure are: $\delta$ and~$\tau_0$; the Haar measure of the skew elements that gives mass $1$
to the integral ones; the modulus \m{\modT{\sT}(u)=\gix{u}^{-1}}; the character \m{\charT{S}{\sT}}; and the
Plancherel constant~$1$. The character is self-dual: for \m{a\in A}, \m{\charT{S}{\sT}(ra)=1} for all \m{r\in O_{\sT}}
exactly if \m{\psiS{S}(r'\tilde a)=1} for all \m{r'\in\hO}, that is, exactly if \m{\tilde a\in\hO}
(Input~\ref{i:1}(b)), that is, \m{a\in O_{\sT}} (Input~\ref{i:5}(a)). The formal text verifies the further axioms
of a place structure for these data: the decomposition of the Haar measure of $A$ along the map
\m{(\beta,\beta')\mapsto\tau_0\delta\beta+\beta'} from pairs of skew elements, the continuity of the modulus, a locally
integrable majorant of \m{\gix{x}^{1/2}}, and Plancherel's inequality, whose kernel is the indicator function of
$O_{\sT}$ by the self-duality.

\emph{The level.} By Input~\ref{i:1}(a), \m{\Tlat{S}{\ET{\sT}}} is the lattice form for~\m{\NS{S}}, and its split
Gram matrix is a hermitian unit. Put \m{\iota=\{1,\dots,\kk{\fs}-6\varrho-13\}} and \m{m'=3\varrho+6}; then
\m{\kk{\fs}=\#\iota+2m'+1}, as \m{6\varrho+13\le\kk{\fs}}. The unit \m{g=\gid{S}_{\sT}} is integral and fixed by the
involution. By Input~\ref{i:4}(c) there is a lattice level $\cL$ with \m{\gL{\cL}=g}. It is standard, with
\m{m=3\varrho+5}: by Input~\ref{i:5}(c), for \m{u=\nn{\fs}}, \m{v=2\varrho} and \m{n=\nn{\fs}}, and by
\m{\nn{\fs}[E:\Q]\le2^{2\varrho}} (Setting~\ref{s:S}), \m{\zeta_A(m)^{\nn{\fs}}\le\zeta_{\hO}(m)^{\nn{\fs}}\le1+2^{-\nn{\fs}}};
\m{m<9\nn{\fs}}; and \m{13\nn{\fs}\le\kk{\fs}-6\varrho-13} because \m{1000(2\varrho+2)\le\kk{\fs}}. Finally the unitary
group \m{\GT{\Tlat{S}{\ET{\sT}}}} has a Haar measure invariant on both sides: by Input~\ref{i:5}(d) the lattice has
an orthonormal basis over~\m{\hO}, which identifies \m{\GFin{S}} with the finite adelic unitary group of the
standard form, a group with a Haar measure invariant on both sides; the elements that are integral off $\sT$
form an open subgroup, on which \m{\loc{S}{\sT}} is surjective by Lemma~\ref{l:lift}; and the image of the
restricted measure is a Haar measure of the image. So \m{\isFL{S}{\sT}{\cL}}.
\end{proof}

\begin{proposition}[counting reductions by a measure]\label{p:count}
In Setting~\ref{s:T}, suppose that \m{\gid{S}^{-1}} is integral off~$\sT$. Let $B$ be a subgroup of \m{\GFin{S}} and
$\mu$ an s-finite measure on \m{\GT{\Tlat{S}{\ET{\sT}}}} such that $\mu$-almost every $z$ is \m{\loc{S}{\sT}(x)} for
some \m{x\in B} with \m{\intoff{S}{\sT}{x}}, and such that $\mu$ is invariant under right translation by
\m{\loc{S}{\sT}(x)} for every \m{x\in B}. Then
\m{\#\redim{S}{B}\cdot\mu\bigl(\cgp{\cL}\bigr)=\mu\bigl(\KT{S}{\sT}{1}\bigr)}.
\end{proposition}
\begin{proof}
For every \m{s\in\redim{S}{B}} choose \m{k_s\in B\cap\KS{S}} with \m{\redS{S}(k_s)=s}, and let
$C_s=\{z:\ z\,\loc{S}{\sT}(k_s)\in\cgp{\cL}\}$. By the invariance, \m{\mu(C_s)=\mu(\cgp{\cL})}. The sets $C_s$ are
pairwise disjoint: if \m{z\in C_s\cap C_t}, then $\loc{S}{\sT}(k_s)^{-1}\loc{S}{\sT}(k_t)\in\cgp{\cL}$, so \m{s=t} by
Lemma~\ref{l:levels}(d). Their union is \m{\KT{S}{\sT}{1}} up to a $\mu$-null set. Indeed, each $C_s$ lies in
\m{\KT{S}{\sT}{1}}, because \m{\cgp{\cL}} and \m{\loc{S}{\sT}(k_s)} do (Lemma~\ref{l:levels}(b)). Conversely let
$z=\loc{S}{\sT}(x)\in\KT{S}{\sT}{1}$ with \m{x\in B}, \m{\intoff{S}{\sT}{x}}. Then \m{x\in\KS{S}} by
Lemma~\ref{l:levels}(b), \m{s=\redS{S}(x^{-1})} lies in \m{\redim{S}{B}}, and \m{xk_s\in\KgS{S}}, so
$z\,\loc{S}{\sT}(k_s)\in\cgp{\cL}$ by Lemma~\ref{l:levels}(c). Summing the measures gives the claim.
\end{proof}
\leanline{\lean{cardreducedImagemul}}

\begin{corollary}[the reduction of the Bessel subgroup]\label{c:bessel}
In Setting~\ref{s:S}, suppose that $E$ carries Weil data, and write \m{\mathcal S=\ST{\Tlat{S}{\AEf}}},
\m{\mathcal H=\HT{\Tlat{S}{\AEf}}}. Then \m{\#\redim{S}{\mathcal S}=\#\redim{S}{\mathcal H}\cdot\uord{S}}.
\end{corollary}
\begin{proof}
By Theorem~C(ii), for \m{\sT_0=\emptyset}, and Theorem~C(i) there are a set $\sT$ off which \m{\gid{S}^{-1}} is
integral and a filling level $\cL$ for $S$ at~$\sT$. Apply Proposition~\ref{p:count} to
\m{B=\mathcal U=\UT{\Tlat{S}{\AEf}}} and \m{\mu=\etab{\varrho}}: this measure is carried by the unipotent group at~$\sT$,
every element of which is the image of an element of $\mathcal U$ that is integral off $\sT$
(Lemma~\ref{l:lift}), and it is invariant under right translation by that group (Input~\ref{i:4}(b)). By the
same input \m{\mu(\cgp{\cL})=\umass{\cL}} and \m{\mu(\KT{S}{\sT}{1})=1}. Hence \m{\#\redim{S}{\mathcal U}\cdot\umass{\cL}=1},
and \m{\#\redim{S}{\mathcal U}=\uord{S}} by Lemma~\ref{l:mod}(c). It remains to see that
\m{(\bar u,\bar h)\mapsto\bar u\bar h} is a bijection from \m{\redim{S}{\mathcal U}\times\redim{S}{\mathcal H}} to
\m{\redim{S}{\mathcal S}}. It is injective because the two images meet in~$1$: if \m{u\in\mathcal U\cap\KS{S}} and
\m{h\in\mathcal H\cap\KS{S}} have the same reduction, then \m{u^{-1}h\in\KgS{S}}, so \m{h\in\KgS{S}} by
Input~\ref{i:1}(e). It is surjective because an element of \m{\mathcal S\cap\KS{S}} is a product $uh$ with
\m{u\in\mathcal U}, \m{h\in\mathcal H}, and then \m{h\in\KS{S}} by the same input, hence also \m{u\in\KS{S}}.
\end{proof}
\leanline{\lean{cardreducedImagebesselS}}

\sect{The floor group, Weyl elements, and blocks as limits}{s:floor}
\begin{lemma}[positions of the elements of the floor group]\label{l:floor}
In Setting~\ref{s:S}, let \m{\gamma_1\in\GS{S}{E}} with \m{\gamma_1\in\Hl{\varrho}{\Trat{S}{E}}} and let $c$, $c'$ be classes.
Put \m{l=\pos{S}{c}{c'}{\gamma_1}}, \m{t=\Zf{S}{\gamma_1}_{\varrho\varrho}}, and let \m{\omega\in E^{\kappa}} be the vector of the
coefficients of \m{\gamma_1z} at the basis vectors of~$\kappa$. Then:
\begin{enumerate}[label=\textup{(\roman*)},leftmargin=2.3em,itemsep=1pt]
\item \m{l\in\Hl{\varrho}{\Tlat{S}{\AEf}}};
\item \m{l\,z=\tailv{\piS{S}}{t}{y_c^{-1}\omega}}, with $t$ and $\omega$ read in \m{\AEf};
\item \m{\NS{S}\cdot\omega^*\omega=t\bar t-\NS{S}^2};
\item for every sequence $p$ of units of~$E$, \m{\pos{S}{c}{c'}{\weyl{\varrho}{p}\gamma_1}=\weyl{\varrho}{p}\cdot l}, where on
the right $p$ is read as a sequence of finite ideles.
\end{enumerate}
\end{lemma}
\begin{proof}
As $\gamma_1$ fixes all $e_a$ and $e^*_a$, it preserves the span of $\kappa$ and~$z$; write
$\gamma_1z=\sum_j\omega_jx_j+\zeta z$. The form is the identity on $\kappa$ and \m{-\NS{S}} at~$z$, so
\m{t=h(z,\gamma_1z)=-\NS{S}\zeta}, and \m{h(\gamma_1z,\gamma_1z)=h(z,z)} reads
$\omega^*\omega-\NS{S}\bar\zeta\zeta=-\NS{S}$; this gives~(iii). (i) The elements \m{\hc{S}{c}}, \m{\hc{S}{c'}} lie in
$\mathcal H$ and \m{\rfin{S}(\gamma_1)} fixes all $e_a$, $e^*_a$, since the change to the lattice basis only
concerns the middle. (ii) \m{\hc{S}{c'}} fixes~$z$, so $l\,z=\hc{S}{c}^{-1}\rfin{S}(\gamma_1)z$; in the rational
basis this is the vector with the $\kappa$-part \m{y_c^{-1}\omega} and the coefficient
$\zeta=-(\piS{S}\overline{\piS{S}})^{-1}t$ at~$z$, and written in the lattice basis it is the tail vector.
(iv) The Weyl element only moves the vectors $e_a$, $e^*_a$; it commutes with the change of basis and with
the elements that fix all $e_a$ and~$e^*_a$, such as \m{\hc{S}{c}} and \m{\hc{S}{c'}}.
\end{proof}
\leanline{\lean{termPositionfloorGroupspec}}

\begin{lemma}[Weyl elements at infinity; tail vectors at $\sT$]\label{l:weylinf}
In Setting~\ref{s:S}:
\begin{enumerate}[label=\textup{(\alph*)},leftmargin=2.3em,itemsep=1pt]
\item For \m{\gamma\in\GS{S}{E}}, a sequence $p$ of units of $E$ and an infinite place $\tau$ of~$E$,
\m{\rinf{S}(\weyl{\varrho}{p}\gamma)_\tau=\weyl{\varrho}{p^{(\tau)}}\cdot\rinf{S}(\gamma)_\tau}, where \m{p^{(\tau)}} is the sequence
\m{a\mapsto\fscale{S}{\tau}{a}\cdot(p_a)_\tau}, read as $1$ at the indices where this number is not a unit.
\item For a finite set $\sT$ of finite places of $E$ that is stable under complex conjugation, a finite
adele $t$ and \m{x\in\AEf^{\kappa}}: \m{\tailv{\piS{S}_{\sT}}{t_{\sT}}{x_{\sT}}=\bigl(\tailv{\piS{S}}{t}{x}\bigr)_{\sT}}.
\end{enumerate}
\end{lemma}
\begin{proof}
(a) \m{\rinf{S}} is a homomorphism, and at the place $\tau$ it conjugates the image of a matrix by the diagonal
matrix of the basis scales. The Weyl element has the entries $p_a$ and \m{\bar p_a^{-1}} between $e_a$ and
$e^*_a$; conjugation multiplies $p_a$ by the quotient of the two scales, which is the frame scale
\m{\fscale{S}{\tau}{a}}, a positive real number, and \m{\bar p_a^{-1}} by its inverse. (b) Taking components at $\sT$
is a homomorphism of rings with involution, and the tail vector is given by ring operations.
\end{proof}
\leanline{\lean{ratToInfweylElemmul}\sep \lean{tailVectorevalPlaces}}

\begin{proof}[Proof of Theorem B]
Write \m{T=\Trat{S}{E}} and \m{\mathcal U=\UT{T}}, and let
\[\mathcal B=\tlink{S}{\Lambda}\cap\terms{S}{c}{s}{c'}.\]
As the amplitude sum over \m{\tlink{S}{\Lambda}} is not zero, there is an element $\gamma_0$ of~$\mathcal B$. By hypothesis its leading
minors do not vanish, so by Input~\ref{i:3}(a) we can write $\gamma_0=v^{-1}\,\weyl{\varrho}{r}\gamma_1\,v'$ with
$v,v'\in\mathcal U$, a sequence $r$ of units and $\gamma_1\in\Hl{\varrho}{T}$. Let $t=h(z,\gamma_1z)=\Zf{S}{\gamma_1}_{\varrho\varrho}$. By Input~\ref{i:3}(c) the last
leading minor of \m{\Zf{S}{\gamma_0}} is \m{r_0\cdots r_{\varrho-1}t}, so $t$ is a unit; the term discriminant of $\gamma_0$
is \m{d\bar d(t\bar t-\NS{S}^2)\ne0}; and $\Lambda=\link{\gamma_0}=\link{\weyl{\varrho}{p}\gamma_1}$ for every sequence~$p$.

\emph{(i)} By Lemma~\ref{l:floor}(iii), at an infinite place~$\tau$,
\m{|t_\tau|^2-\NS{S}^2=\NS{S}\sum_j|(\omega_j)_\tau|^2\ge0}, and it is not zero because \m{t\bar t\ne\NS{S}^2} in~$E$.
The element \m{k=\pos{S}{c}{c'}{\gamma_0}} lies in \m{\KS{S}} and has the reduction~$s$, as \m{\gamma_0\in\terms{S}{c}{s}{c'}}.

\emph{(ii)} For \m{p\colon\{0,\dots,\varrho-1\}\to E^{\times}} and \m{u,u'\in\mathcal U} put
\m{\gamma_{p,u,u'}=u\,\weyl{\varrho}{\pad{p}{t}}\gamma_1\,u'^{-1}}. The map \m{(p,u,u')\mapsto\gamma_{p,u,u'}} is injective by Input~\ref{i:3}(b). Its values lie
in \m{\tlink{S}{\Lambda}}: the link is $\Lambda$, and the term discriminant is \m{d\bar d(t\bar t-\NS{S}^2)\ne0} with
\m{d=p_0\cdots p_{\varrho-1}} (Input~\ref{i:3}(c)). Every \m{\gamma\in\mathcal B} is a value: by Input~\ref{i:3}(a),
\m{\gamma=w^{-1}\weyl{\varrho}{q}\gamma_2w'} with \m{\gamma_2\in\Hl{\varrho}{T}}, and
\m{\link{\weyl{\varrho}{q}\gamma_2}=\Lambda=\link{\weyl{\varrho}{1}\gamma_1}} forces \m{\gamma_2=\gamma_1}. Therefore the family
\m{\mathbf 1[\gamma_{p,u,u'}\in\terms{S}{c}{s}{c'}]\cdot\amp{\sigma}{S}{c}{c'}{\gamma_{p,u,u'}}}, indexed by the triples, is a
reindexing of the summable family of the amplitudes over $\mathcal B$ (Input~\ref{i:2}(c)), completed by zeros.
Its sum is the amplitude sum of the theorem, and summing first over \m{(u,u')} gives a summable family
\[L_p=\sum_{(u,u')}\mathbf 1\bigl[\gamma_{p,u,u'}\in\terms{S}{c}{s}{c'}\bigr]\cdot\amp{\sigma}{S}{c}{c'}{\gamma_{p,u,u'}}\]
with the same sum. The coset amplitude \m{\cosamp{S}{\sT}{c}{s}{c'}{\weyl{\varrho}{\pad{p}{t}}\gamma_1}} is the same sum
with the further factor \m{\mathbf 1[\intoff{S}{\sT}{\pos{S}{c}{c}{u}}]}. For each \m{(u,u')} this factor is $1$ for
all sufficiently large~$\sT$ (Theorem~C(iii)), and the terms are dominated by the norms of a summable family;
by dominated convergence for series the coset amplitudes tend to~$L_p$.
\end{proof}

\sect{The unfolding}{s:unfold}
\begin{proof}[Proof of Theorem A]
Write $G_{\sT}=\GT{\Tlat{S}{\ET{\sT}}}$, $g=\gL{\cL}$, $\mathcal U=\UT{\Trat{S}{E}}$, $s=\redS{S}(k)$ and
$P=\pos{S}{c}{c'}{\gamma}$, which is integral off~$\sT$; for $u,u'\in\mathcal U$ put $\gamma_{u,u'}=u\gamma u'^{-1}$.

\emph{Step 1: the left side.} By Lemma~\ref{l:floor}(iv), $\loc{S}{\sT}(P)=\weyl{\varrho}{p_{\sT}}\cdot l_0$ with
$l_0=\loc{S}{\sT}(\pos{S}{c}{c'}{\gamma_1})$, so that
$\corr{Z}{p_{\sT}}{l_0^{-1}}=\hdens{\cL}{\varrho}{p_{\sT}}{1}\cdot\bcorr{\gM{\cL}}{\bvec{Z}{\cdot}}{\loc{S}{\sT}(P)}$. Its coordinate at
\m{(a,b)} is the head density times
\begin{align*}&\int\phs{g}{z\,a_a^{-1}}\,\overline{\phs{g}{\loc{S}{\sT}(P)^{-1}z\,a_b^{-1}}}\,d\gM{\cL}(z)\\
&\qquad=\int\phs{g}{z}\,\overline{\phs{g}{\loc{S}{\sT}(P)^{-1}z\,\loc{S}{\sT}(k)}}\,d\gM{\cL}(z),\end{align*}
by the right invariance of \m{\gM{\cL}} (substitute \m{z\,a_a} for~$z$), the hypothesis
$a_aa_b^{-1}\in\loc{S}{\sT}(k)\,\cgp{\cL}$ and Input~\ref{i:4}(a). The archimedean factor is
$\bcorr{dy}{W^{\sigma}}{\rinf{S}(\gamma)}=\int\bker{\sigma}{S}{\gamma}{y}\,dy$. So the left side of the theorem is
$\hdens{\cL}{\varrho}{p_{\sT}}{1}$ times the integral over \m{(z,y)} of
\begin{equation}\label{e:r}\mathsf r(z,y)=\phs{g}{z}\cdot\overline{\phs{g}{\loc{S}{\sT}(P)^{-1}z\,\loc{S}{\sT}(k)}}\cdot\bker{\sigma}{S}{\gamma}{y}.\end{equation}

\emph{Step 2: the right side.} For \m{(u,u')} let \m{Q(u,u')} be the condition of Definition~\ref{d:cosamp}:
\m{\intoff{S}{\sT}{\pos{S}{c}{c}{u}}} and \m{\gamma_{u,u'}\in\terms{S}{c}{s}{c'}}. By the left invariance of \m{\gM{\cL}} and of
$dy$, for every \m{(u,u')},
\begin{align*}&\gM{\cL}\bigl(\cgp{\cL}\bigr)\cdot\amp{\sigma}{S}{c}{c'}{\gamma_{u,u'}}\\
&\qquad=\int\!\!\int\mathbf 1\bigl[\loc{S}{\sT}(\pos{S}{c}{c}{u})\,z\in\cgp{\cL}\bigr]
\cdot\mathsf a_{u,u'}\bigl(\rinf{S}(u)\,y\bigr)\,d\gM{\cL}(z)\,dy,\end{align*}
where \m{\mathsf a_{u,u'}(y)=\cutoff{S}{c}(y)\,\cutoff{S}{c'}(\rinf{S}(\gamma_{u,u'})^{-1}y)\,\bker{\sigma}{S}{\gamma_{u,u'}}{y}} is the
integrand of the amplitude. The sum over the \m{(u,u')} with \m{Q(u,u')} may be taken under the integral: the
integral of the norm of the term of \m{(u,u')} is at most
\m{\gM{\cL}(\KT{S}{\sT}{1})\cdot\ampabs{\sigma}{S}{c}{c'}{\gamma_{u,u'}}}, the map \m{(u,u')\mapsto\gamma_{u,u'}} is injective (Input~\ref{i:3}(b)), and
the sum of these numbers over \m{\terms{S}{c}{s}{c'}} is finite (Input~\ref{i:2}(c); the measure is finite by
Input~\ref{i:4}(b)). So it suffices to show, for all $z$ and~$y$,
\begin{equation}\label{e:pt}\sum_{(u,u')}\mathbf 1\bigl[Q(u,u')\bigr]\cdot\mathbf 1\bigl[\loc{S}{\sT}(\pos{S}{c}{c}{u})\,z\in\cgp{\cL}\bigr]\cdot
\mathsf a_{u,u'}\bigl(\rinf{S}(u)\,y\bigr)=\mathsf r(z,y).\end{equation}

\emph{Step 3: the two conditions.} Fix $z$ and~$y$. By Lemma~\ref{l:lift}, \m{z=\loc{S}{\sT}(x)} with
\m{\intoff{S}{\sT}{x}}. Put \m{\xi=P^{-1}xk}; it is integral off~$\sT$ (Lemma~\ref{l:levels}(b) for~$k$), and
\m{\loc{S}{\sT}(\xi)=\loc{S}{\sT}(P)^{-1}z\,\loc{S}{\sT}(k)}. We claim that, for \m{u,u'\in\mathcal U},
\begin{equation}\label{e:C}\begin{gathered}Q(u,u')\text{ and }\loc{S}{\sT}(\pos{S}{c}{c}{u})\,z\in\cgp{\cL}\\
\Longleftrightarrow\quad\pos{S}{c}{c}{u}\,x\in\KgS{S}\text{ and }\pos{S}{c'}{c'}{u'}\,\xi\in\KgS{S}.\end{gathered}\end{equation}
By Lemma~\ref{l:levels}(c), and since $x$ is integral off~$\sT$, the element \m{l_1=\pos{S}{c}{c}{u}\,x} lies in
\m{\KgS{S}} exactly if \m{\pos{S}{c}{c}{u}} is integral off $\sT$ and \m{\loc{S}{\sT}(\pos{S}{c}{c}{u})\,z\in\cgp{\cL}}. Suppose
this. From \m{P=xk\xi^{-1}} we get
\[\pos{S}{c}{c'}{\gamma_{u,u'}}=\pos{S}{c}{c}{u}\cdot P\cdot\pos{S}{c'}{c'}{u'}^{-1}=l_1\cdot k\cdot\bigl(\pos{S}{c'}{c'}{u'}\,\xi\bigr)^{-1}.\]
As \m{l_1\in\KgS{S}} and \m{k\in\KS{S}} has the reduction~$s$, the left side lies in \m{\KS{S}} with the reduction $s$
exactly if \m{\pos{S}{c'}{c'}{u'}\,\xi\in\KgS{S}}. This proves~\eqref{e:C}.

\emph{Step 4: phases.} Let \m{u\in\mathcal U} with \m{\pos{S}{c}{c}{u}\,x\in\KgS{S}}. Then
\begin{equation}\label{e:val}\phs{g}{z}=\uchar{S}{\rfin{S}(u)}^{-1}.\end{equation}
Indeed \m{u_1=\pos{S}{c}{c}{u}} lies in \m{\UT{\Tlat{S}{\AEf}}} and is integral off $\sT$ (Input~\ref{i:1}(c), Step~3); by
Lemma~\ref{l:levels}(c), \m{\loc{S}{\sT}(u_1)\,z\in\cgp{\cL}}, so by Input~\ref{i:4}(a) and Lemma~\ref{l:char}
\m{\phs{g}{z}=\psiT{\sT}\bigl(g^{-1}\sd{\loc{S}{\sT}(u_1)^{-1}}\bigr)=\uchar{S}{u_1}^{-1}}; and \m{\uchar{S}{u_1}=\uchar{S}{\rfin{S}(u)}} by
Input~\ref{i:1}(c), because \m{u_1=\hc{S}{c}^{-1}\rfin{S}(u)\hc{S}{c}}. The same holds for $u'$, $c'$ and~$\xi$. With
Input~\ref{i:2}(b) and the definition of the cutoffs' arguments, for \m{(u,u')} satisfying \eqref{e:C} the term
of \eqref{e:pt} is
\begin{align*}&\cutoff{S}{c}\bigl(\rinf{S}(u)y\bigr)\cdot\cutoff{S}{c'}\bigl(\rinf{S}(u')\rinf{S}(\gamma)^{-1}y\bigr)\cdot\uchar{S}{\rfin{S}(u)}^{-1}
\uchar{S}{\rfin{S}(u')}\cdot\bker{\sigma}{S}{\gamma}{y}\\
&\qquad=\cutoff{S}{c}\bigl(\rinf{S}(u)y\bigr)\cdot\cutoff{S}{c'}\bigl(\rinf{S}(u')\rinf{S}(\gamma)^{-1}y\bigr)\cdot\mathsf r(z,y),\end{align*}
by \eqref{e:val} for \m{(u,x)} and for \m{(u',\xi)}, the values of $\uchar{S}{\cdot}$ being of modulus~$1$.

\emph{Step 5: the cutoffs sum to one.} So the left side of \eqref{e:pt} is \m{\mathsf r(z,y)} times the product of
\[\Sigma_c(x,y)=\sum_{u\in\mathcal U}\mathbf 1\bigl[\pos{S}{c}{c}{u}\,x\in\KgS{S}\bigr]\cdot\cutoff{S}{c}\bigl(\rinf{S}(u)\,y\bigr)\]
and of \m{\Sigma_{c'}(\xi,\rinf{S}(\gamma)^{-1}y)}. If \m{\mathsf r(z,y)=0} we are done. Otherwise \m{\phs{g}{z}\ne0} and
\m{\phs{g}{\loc{S}{\sT}(\xi)}\ne0}, and it suffices to show: if $x$ is integral off $\sT$ and
\m{\phs{g}{\loc{S}{\sT}(x)}\ne0}, then \m{\Sigma_c(x,y)=1} for all~$y$. By Input~\ref{i:4}(a) and Lemma~\ref{l:lift} there
is \m{u_2\in\UT{\Tlat{S}{\AEf}}}, integral off~$\sT$, with \m{\loc{S}{\sT}(u_2^{-1}x)\in\cgp{\cL}}, hence
\m{u_2^{-1}x\in\KgS{S}} (Lemma~\ref{l:levels}(c)). By Input~\ref{i:1}(d) there is a rational \m{v\in\mathcal U} with
\m{u_2^{-1}\pos{S}{c}{c}{v}\in\KgS{S}}. Then, for \m{u\in\mathcal U},
\m{\pos{S}{c}{c}{u}\,x\in\KgS{S}} exactly if \m{\pos{S}{c}{c}{uv}\in\KgS{S}}, that is, \m{uv\in\clU{S}{c}}; so with \m{u=wv^{-1}},
\[\Sigma_c(x,y)=\sum_{w\in\clU{S}{c}}\cutoff{S}{c}\bigl(\rinf{S}(w)\,\rinf{S}(v)^{-1}y\bigr)=1\]
by the partition of unity of Input~\ref{i:2}(a). This proves \eqref{e:pt} and the theorem.
\end{proof}

\sect{The filling data of the admissible data}{s:data}
\begin{definition}[the filling data of $D$]\label{d:SD}
In Setting~\ref{s:D}, the \emph{filling data of~$D$}, \m{\SD{D}}, are the filling data over \m{E(D)} of the shape
\m{\fs(D)} with the following components: the number \m{b(D)}; the prime \m{N(D)} and the idele \m{\hat\pi(D)}; the
primes \m{g_0(D)}, \m{g_1(D)} with the chosen places of $D$ above them; the generator \m{\varpi(P)} of the
different; the column level \m{\Colv{D}{0}}; and at every infinite place $\tau$ a Bessel system chosen among
those of Proposition~\ref{p:bessel}. The conditions on filling data (Setting~\ref{s:S}) hold by
Input~\ref{i:6}(a),~(b). In particular \m{\bS{\SD{D}}=b(D)}, \m{\NS{\SD{D}}=N(D)}, \m{\gS{\SD{D}}=g_0(D)g_1(D)}, and for
the shape \m{\fs=\fs(D)}: \m{\rh{\fs}=\varrho(D)}, \m{\kk{\fs}=k(D)} and \m{\rr{\fs}=2\varrho(D)+k(D)+1=r(D)}.
\end{definition}
\leanline{\lean{filling}}

\begin{lemma}[three small facts]\label{l:small}
In Setting~\ref{s:D}: the index set \m{X_0(D)} of the columns has \m{k(D)} elements. Moreover the formal text
fixes, as instances, a unitary matrix that diagonalises the real form of the shape \m{\fs(D)} with the signs of
the space of the rows (a \emph{frame}), and records the bound \m{b(D)^{10}\le2^{\varrho(D)}}.
\end{lemma}
\begin{proof}
The first is the count of the index set in the unit on the seesaw. The frame exists because the form has as
many negative signs as the standard form of these signs, and one is chosen. The bound is one of the
logarithmic bounds of Chapter~III: \m{10\log b(D)\le\varrho(D)\log2}.
\end{proof}
\leanline{\lean{cardX}\sep \lean{unitarySignFrameY}\sep \lean{instLbHPowRealNatOfNatHMulLb}}

\begin{proposition}[Bessel systems for the data]\label{p:bessel}
In Setting~\ref{s:D}, let $\tau$ be an infinite place of \m{E(D)}. Then there is a Bessel system $\mathfrak B$ on the
real group \m{G_1(\fs(D))}, for the bound \m{b(D)} and the number \m{f(D)} in the sense of the set-up, of the type
\m{\IsBes{\mathfrak B}{2\varkappa_D(\tau,\cdot)}{\mathsf{w}^{\mathrm{col}}_D(\tau)}}, which means:
\begin{enumerate}[label=\textup{(\roman*)},leftmargin=2.3em,itemsep=1pt]
\item the degree of $\mathfrak B$ is the value of Weyl's dimension polynomial at the doubled row parameter
\m{2\varkappa_D(\tau,\cdot)};
\item there is an invertible matrix $A$ such that, for every $h$ in the unitary group of the columns, embedded
as the subgroup that fixes the frame and its dual, and every point~$y$, the full Bessel function satisfies
\m{A\,W^{\full}(hy)=\mathrm{Sch}(h)\,A\,W^{\full}(y)}, where \m{\mathrm{Sch}} is the transpose of the Schur
representation of weight \m{\mathsf{w}^{\mathrm{col}}_D(\tau)} at the transposed element; and
\item $W^{\full}$ is fixed by right convolution with the kernel of the discrete series of the parameter
\m{2\varkappa_D(\tau,\cdot)}, read in the frame of Lemma~\ref{l:small}.
\end{enumerate}
\end{proposition}
\begin{proof}
An assembly of Input~\ref{i:6}(d); its content lies there. The steps, in the order of the formal text: (1) the
doubled row parameter $\lambda$ has entries of absolute value at most \m{3a(D)} and a gap of size \m{r(D)^7}
(Chapter~III), and its lowest type for the maximal compact subgroup is the column weight at one position and
$0$ at the others. (2) The formal degree of $\lambda$ is positive and equals Weyl's dimension polynomial; the
kernel of its discrete series in the frame is idempotent; and the parameter is steep in the sense of the unit
on the exchange. (3) The Gram matrix of the Weyl module of the column weight is positive definite, of the form
$B^*B$; conjugating the Schur representation by $B$ gives a unitary representation~$\rho$, and a matrix that
commutes with it is scalar. (4) There is a Bessel function of type $\rho$ whose matrix coefficient is that of
the lowest type, with the bounds of the Laplace method. (5) It has a cut pair, whose full function satisfies
the Schur bounds. (6) The cut function has band limits; the constant that this requires is bounded as follows:
each factor \m{|\lambda_i-\lambda_j|/2} of Weyl's polynomial is at most \m{3a(D)}, since both entries are at most
\m{3a(D)} in absolute value, and the remaining factorials and powers are compared with powers of \m{b(D)} by the
bounds of Chapter~III. These data form a Bessel system; (i) holds by~(2), (ii) with \m{A=B^*} by~(3), and (iii)
by the convolution identity of the cut pair and the idempotent kernel.
\end{proof}
\leanline{\lean{existsbesselSystem}}

\begin{proof}[Proof of Theorem D]
Write \m{S=\SD{D}}, \m{N=N(D)}, \m{k=k(D)}, \m{f=f(D)}.

\emph{(i)} By Input~\ref{i:6}(c), \m{\incH{S}^{-1}(\KgS{S})} is the group of the \m{y\in K_f} whose row at $k_0$ is
congruent to that of the identity modulo every place dividing \m{\hat\pi(D)}, and whose reduction at every
chosen place is~$1$. By Input~\ref{i:6}(b) this is the description of \m{\Colv{D}{0}}.

\emph{(ii)} By Input~\ref{i:6}(c), \m{\incH{S}^{-1}(\KS{S})\cap K_f} is the group of the \m{y\in K_f} such that, for
each of the $f$ places $w$ dividing \m{\hat\pi(D)}, the inverse transpose of the reduction of $y$ at $w$ fixes the
basis vector of $k_0$ in \m{\F_w^{k}}. No two of these places are conjugate, so the reduction maps $K_f$ onto the
product of the groups \m{\GL_k(\F_w)}, and the index is the product of the indices of the stabilisers of a
non-zero vector, each equal to \m{N^k-1} (Input~\ref{i:5}(e)).

\emph{(iii)} The class automorphism groups of the column level are trivial, so \m{\#\Cl{S}} is its class sum,
which is \m{[K_f:\Colv{D}{0}]\cdot\mass{E(D)}{k}} (Input~\ref{i:6}(b)). By (i) the index is
\[[K_f:\incH{S}^{-1}(\KS{S})\cap K_f]\cdot[\incH{S}^{-1}(\KS{S}):\incH{S}^{-1}(\KgS{S})]=(N^k-1)^f\cdot\#\redim{S}{\mathcal H},\]
by (ii), and because \m{\incH{S}} maps onto $\mathcal H$, so that the second index is the order of the image of
\m{\mathcal H\cap\KS{S}} under~\m{\redS{S}}.

\emph{(iv)} The standing hypotheses are: (1) $\KcolS{S}=\incH{S}^{-1}(\KgS{S})$; (2) \m{\redS{S}} is surjective;
(3) \m{|\vpS{S}_\tau|\le\bS{S}} and $1\le|\kapS{S}{\tau}|\le\bS{S}^{23}$ for all~$\tau$, and
\m{\NS{S}\le\bS{S}^{60\nn{\fs}}}; (4) the band inequality
$\prod_\tau c_{\mathrm{band}}(\fs)\,\bS{S}^{7\kk{\fs}}\bigl(|\kapS{S}{\tau}|^{2\rh{\fs}+1}/(\gS{S}\NS{S})\bigr)^{\kk{\fs}-\nn{\fs}^7}\le\bdeg{S}^{-1}$;
(5) $\tailS{S}^{44}\besE{S}\le\mainE{S}$, $\tailS{S}^{18}\codimF{S}\mainE{S}\le\besE{S}$ and \m{\codimF{S}\le1}; (6)
$\#\Bc{S}\le\bS{S}^{2\kk{\fs}\ninf{E}}$; (7) $(\NS{S}^{\kk{\fs}}/\bS{S})^{\ninf{E}}$ is at most the index of~(ii); and
(8) the bound for the packet constant \m{\packC{S}} stated in Note~F. Here \m{\kapS{S}{\tau}}, \m{c_{\mathrm{band}}(\fs)},
\m{\bdeg{S}}, \m{\tailS{S}}, \m{\besE{S}}, \m{\mainE{S}}, \m{\codimF{S}} are the constants of the set-up, as recalled in
Note~F. (1) is~(i), and (2) is Input~\ref{i:6}(c). The others are inequalities between positive real numbers.
The formal proof makes every quantity explicit: $\#\redim{S}{\mathcal S}=\#\redim{S}{\mathcal H}\cdot\uord{S}$
(Corollary~\ref{c:bessel}), \m{\#\Cl{S}} by~(iii), the index by~(ii), \m{\prod_\tau|\vpS{S}_\tau|} as a power of the
root discriminant, the volumes and the covolume by their definitions, and \m{\rr{\fs}=k+2\varrho(D)+1}. It then
takes logarithms; with the bounds for \m{\#\Bc{S}}, for \m{\log\bdeg{S}} and for the ratios of the masses of the
ranks from $k$ to \m{r(D)} (Input~\ref{i:6}(e)), with \m{N\le N^k} and \m{k\log N-2/N\le\log(N^k-1)}, each inequality
becomes a linear inequality between the logarithmic sizes of the data, which follows from the windows of
Chapter~III (Input~\ref{i:6}(a)). This linear arithmetic is not reproduced here.
\end{proof}

\begin{lemma}[an expansion of a product of stacked determinants]\label{l:poly}
In Setting~\ref{s:D}, with the notation of the note \emph{The height of a ratio of lift volumes}: let $j$ be a
lift index, $i$ a line choice with its short rows~$Y_i$, and let $\Delta$ be the product of the stacked
determinants of the matrices \m{\detRows{j}{i}{Y_i}{\tau}} along \m{\posEquiv{j}}, a polynomial in variables
\m{x_{\tau,a,b}} with $\tau$ an infinite place, \m{a<2m(D)}, and $b$ either a number below \m{k(D)+j-1} or an index of
the space~$L^{+}$. Let $n$ be a set, $f$ a map from the row indices to~$n$ and $g$ a map from the indices of
$L^{+}$ to~$n$. Let $\Delta^{f,g}$ be the image of $\Delta$ under the substitution
\m{x_{\tau,a,b}\mapsto x_{\tau,a,f(\basisVec{D}{b})}} for the numbers $b$ and \m{x_{\tau,a,b}\mapsto x_{\tau,a,g(b)}} for the
indices of~$L^{+}$. Then
\[\Delta^{f,g}=\sum_{l}f_*(P_l)\cdot g_*\bigl(\linePolyCoords{j}{i}{Y_i}_l\bigr),\]
where $l$ runs over the index set of the basis of the Weyl module of the weight of the first columns, $P_l$
is the \emph{letter polynomial} of the $l$-th basis vector for the positive letters (its polynomial, with the
variable of \m{(\tau,a)} and of the number $h$ replaced by the sum of the variables \m{x_{\tau,a,i'}} over the row
indices $i'$ of positive sign and of order~$h$, for the order function of the untwisted auxiliary pair), and $f_*$, $g_*$ rename the last index of the variables.
\end{lemma}
\begin{proof}
By Input~\ref{i:6}(f) and the linearity of the polynomial of a vector, $\Delta$, read as a polynomial in the
first columns with coefficients in the variables of~$L^{+}$, is \m{\sum_lp_l\cdot\linePolyCoords{j}{i}{Y_i}_l} with
$p_l$ the polynomial of the $l$-th basis vector. The substitution of the statement and the map that sends
$p_l$ to \m{f_*(P_l)} and renames the coefficients by $g$ are homomorphisms of algebras, so it suffices to
compare them on the variables. On a variable of~$L^{+}$ both are~$g_*$. On the variable of \m{(\tau,a)} and a
number~$h$, the second gives the sum of the \m{x_{\tau,a,f(i')}} over the row indices $i'$ of positive sign and
order~$h$; there is exactly one such index, \m{i'=\basisVec{D}{h}}, because the order of such an index is its
number, the numbering of the row indices is a bijection, and the sign is positive exactly for the numbers below
\m{k(D)+j-1}.
\end{proof}
\leanline{\lean{stdClassSumcarrierXieqlineMultiplier}}

\sect{Where the results are used}{s:used}
Note~F (\emph{Translates of Poincaré series fill the rows}) quotes Theorem~C as its Input~4(a); Lemma~\ref{l:lift},
Lemma~\ref{l:levels}, Lemma~\ref{l:char} and Lemma~\ref{l:mod} as Input~4(b); Proposition~\ref{p:count} and
Corollary~\ref{c:bessel} as Input~4(c); Lemma~\ref{l:floor} and Lemma~\ref{l:weylinf} as Input~4(d); Theorem~B as
Input~4(e); Theorem~A as Input~4(f); and Definition~\ref{d:SD}, Proposition~\ref{p:bessel} and Theorem~D as
Input~9(a). Lemma~\ref{l:poly} is used in the unit on the rows, in the module on the separation of the lines.
Theorem~D(ii),~(iii) and Corollary~\ref{c:bessel} are used in the proof of Theorem~D(iv).

\sect{Notes on the formal text}{s:notes}
\begin{enumerate}[leftmargin=1.6em,itemsep=3pt]
\item \emph{Arguments that the census reads differently from Lean; to be settled by} \texttt{\#check} \emph{on a
build.} In the three statements of Theorem~C(ii),~(iii) the set $\sT$ is bound by ``for all sufficiently
large'' or ``there are arbitrarily large''; the census lists it as an argument, with its stability under
conjugation, because a variable of the file has this name. The same holds for the stability of $\sT$ in
Theorem~B and Corollary~\ref{c:bessel}, where $\sT$ is bound or absent. For Lemma~\ref{l:mod}(c) the census lists
no hypothesis although the statement names the level; the note states the setting that Lean takes.
\item \emph{Statements cut by the census.} The statements of Lemma~\ref{l:floor} and of Theorem~A contain local
definitions at which the census ends the statement; they were read to their true end, with the second
reader's table of corrected spans.
\item \emph{The side.} The side is an instance argument of the coset amplitude and of Theorems A and~B, with
the full side as default instance; the note writes it as the superscript~$\sigma$.
\item \emph{Forms of the formal statements.} In Theorem~A the two sides are written with scalar
multiplications, and the measure of the congruence subgroup is taken as a real number. In Theorem~B the family
is indexed by the maps from \m{\{0,\dots,\varrho-1\}}, the element $k$ is given as an element of \m{\KS{S}}, and the
limit is along all finite sets of finite places, not only the stable ones. In Definition~\ref{d:FL} the index
set of the orthonormal part and the number of further pairs are fixed by the type of the level. In
Lemma~\ref{l:mod}(b) the right side is written as the square of the inverse of \m{\gS{S}^{\ninf{E}}}.
\item \emph{A hypothesis stronger than its use.} Proposition~\ref{p:count} assumes the invariance of $\mu$ under
the images of all elements of~$B$; the proof uses it for elements of \m{B\cap\KS{S}} only. Note~F states its
Input~4(c) with the invariance under the images of the elements of $B$ that are integral off~$\sT$, which is what
the proof uses and is weaker than the formal hypothesis.
\item \emph{Definition~\ref{d:SD}.} The formal definition gives several components of the filling data only
through the proofs of their properties; they were read off these proofs. The identification of the group
\m{\GRes{\SD{D}}} with \m{\sGD{D}}, used in Note~F, is not a statement of the modules of this note.
\item \emph{Lemma~\ref{l:poly}.} Its formal name speaks of a class sum and a line multiplier; the statement is
the polynomial identity given here, and the class sums appear where it is used.
\item \emph{Points where the dossier's first reading was corrected by the second reader.} In
Lemma~\ref{l:mod}(a) the summary of the proof gave the wrong power of the index; the statement is as above. In
the proof of Theorem~C(i) the self-duality is the statement about the annihilator of~$O_{\sT}$, with the
quantifier over~$r$. In Proposition~\ref{p:bessel} the bound \m{3a(D)} holds for each entry of the parameter
(so \m{6a(D)} for a difference, before the division by~$2$). In Theorem~D(iv) the elementary fact used is
\m{N\le N^k}.
\item \emph{What was not re-derived.} The linear arithmetic of Theorem~D(iv); the constructions behind
Proposition~\ref{p:bessel}; the axioms of a place structure in Theorem~C(i) other than the self-duality; in
Lemma~\ref{l:poly}, that the order of a row index of positive sign is its number, which is taken from the last
step of the formal proof. The measurability and summability conditions of the interchange of sum and integral
in Theorem~A are verified in the formal text; the text above gives the bound that is used.
\end{enumerate}

\subhead{How this note was written.} From the structured proofs of the three modules in the third chunk of the
unit's dossier, all of whose statements and claims a second reader had re-read; her remarks were taken into
account. I read the three modules in the Lean text line by line and compared every statement of this note
with it argument by argument; I read in the Lean text the definitions of the filling data, of admissibility, of
Bessel systems and their type, of the character at~$\sT$, and of the notions of the unit on the rows that
Lemma~\ref{l:poly} needs. The inputs are stated from the dossier's readings of their quoted statements.
Arguments that are mine and shorter than the formal ones: the proof of Lemma~\ref{l:lift} by the identity for
\m{x^*Jx-J}, the proof of Lemma~\ref{l:floor} in coordinates, and the organisation of the proof of Theorem~A in
five steps. Nothing was compiled for this note. The program check of the unit's dossier covered \UnitSteps{}
leaf steps of \UnitProofs{} proofs.

\head{Appendix A. Concordance with the Lean artifact}
{\raggedright The formal names of the objects and the letters of this note: \texttt{IsIntegralOff},
\texttt{toPlaces}, \texttt{IsFillingLevel}, \texttt{cosetAmplitude} are \m{\intoff{S}{\sT}{x}}, \m{\loc{S}{\sT}},
\m{\isFL{S}{\sT}{\cL}}, \m{\cosamp{S}{\sT}{c}{s}{c'}{\gamma}}; \texttt{Data.filling} is \m{\SD{D}}; \texttt{LocProd},
\texttt{evalPlaces}, \texttt{adeleOfLocal}, \texttt{integralOff} are \m{\ET{\sT}}, \m{x\mapsto x_{\sT}},
\m{y\mapsto\tilde y}, integrality off~$\sT$; \texttt{unitaryLevel} is \m{\KA{S}{u}} over the finite adeles and
\m{\KT{S}{\sT}{u}} over \m{\ET{\sT}}; \texttt{latStabilizer}, \texttt{principalLevel}, \texttt{reduction},
\texttt{reducedImage} are \m{\KS{S}}, \m{\KgS{S}}, \m{\redS{S}}, \m{\redim{S}{B}}; \texttt{termPosition},
\texttt{termsOver}, \texttt{termsOfLink} are \m{\pos{S}{c}{c'}{\gamma}}, \m{\terms{S}{c}{s}{c'}}, \m{\tlink{S}{\Lambda}};
\texttt{selfDualChar}, \texttt{charOnPlaces}, \texttt{unipotentChar} are $\psiS{S}$, \m{\charT{S}{\sT}},
\m{\uchar{S}{u}}; \texttt{besselKernel}, \texttt{classCutoff}, \texttt{amplitude}, \texttt{amplitudeSum} are
\m{\bker{\sigma}{S}{\gamma}{y}}, \m{\cutoff{S}{c}}, \m{\amp{\sigma}{S}{c}{c'}{\gamma}}, \m{\ampsum{\sigma}{S}{\Xi}{c}{s}{c'}};
\texttt{unipOrder}, \texttt{unipotentMass}, \texttt{frameScaleSq} are \m{\uord{S}}, \m{\umass{\cL}},
\m{\fscale{S}{\tau}{a}}; \texttt{padTuple}, \texttt{tailVector}, \texttt{weylElem} are \m{\pad{p}{t}},
\m{\tailv{\pi}{t}{x}}, \m{\weyl{\varrho}{p}}; the element \texttt{σ} of the formal statements is $\gamma_1$, the set
\texttt{T} is~$\sT$, the level \texttt{F} is~$\cL$; \texttt{paramW} and \texttt{Data.κ} are the doubled row
parameter and the column weight. Two statements as the artifact has them:\par}
\quoted{correlationsmuleqcosetAmplitude}
\quoted{existshasSumcosetAmplitude}
{\fontsize{8.6}{11.6}\selectfont\renewcommand{\leanfont}{\ttfamily\fontsize{7.7}{10}\selectfont\color{leancol}}
\setlength{\LTleft}{0pt}\setlength{\LTright}{\fill}\setlength{\tabcolsep}{0pt}
\begin{longtable}{@{}p{27mm}>{\raggedright\arraybackslash}p{68mm}p{34mm}r@{}}
\textbf{In this note} & \textbf{Lean declaration} & \textbf{File} & \textbf{Lines}\\\midrule\endhead
\ConcordanceRows
\end{longtable}}

\head{Notation}
\begingroup\fontsize{8.6}{10.9}\selectfont\raggedright
\setlength{\columnsep}{9mm}\setlength{\parskip}{0pt}
\begin{multicols}{2}
\nt{$E$, \m{x\mapsto\bar x}, $\tau$, $v$, $x_\tau$, \m{\ninf{E}}}{a CM field with its CM type, conjugation,
infinite places, the image at a place, their number}
\nt{$w$, $\bar w$, \m{\F_w}, $E_w$, $O_w$, \m{\AEf}, \m{\hO}, $\psi_f$}{a finite place, its conjugate, residue
field, completion, integers; finite adeles, integral adeles, the standard character}
\nt{$\sT$, $\sT_0$, $\sT'$, \m{\ET{\sT}}, $O_{\sT}$, $x_{\sT}$, $\tilde y$}{finite sets of places (of units for
$\sT_0$), the product of completions, its integers, components, extension by zero}
\nt{\m{\psiT{\sT}}, \m{\modT{\sT}}, \m{\gix{u}}, $\delta$, $\tau_0$, $\beta$, $A$}{character, modulus, generalised index, a
skew unit, a trace-one element and skew elements of the place structure on \m{A=\ET{\sT}}}
\nt{\m{\fs=(\varrho,\kappa,k_0)}, \m{\rh{\fs}}, \m{\kk{\fs}}, \m{\nn{\fs}}, \m{\rr{\fs}}}{a shape and its numbers; in
proofs $\varrho$ abbreviates \m{\rh{\fs}}}
\nt{\m{\Ix{\fs}}, $e_a$, $e^*_a$, $z$, $x_j$, \m{\mathfrak f_a}, $h$}{index set, basis vectors, the frame, the
hermitian form}
\nt{$T$, \m{\JT{T}}, \m{\GT{T}}, \m{\ZT{T}{x}}, \m{\UT{T}}, \m{\HT{T}}, \m{\ST{T}}, \m{\Hl{\varrho}{T}}}{a form, split Gram
matrix, unitary group, frame Gram matrix, Bessel unipotent subgroup, the subgroup fixing the frame and its
dual, Bessel subgroup, floor group; $\mathcal U$, $\mathcal H$, $\mathcal S$ in proofs}
\nt{\m{\sd{u}}, \m{\weyl{\varrho}{p}}, \m{\pad{p}{t}}, \m{\tailv{\pi}{t}{x}}}{sum above the diagonal, Weyl element,
padded sequence, tail vector}
\nt{$S$, \m{\bS{S}}, \m{\NS{S}}, \m{\piS{S}}, \m{\qS{S}{i}}, \m{\gS{S}}, \m{\ChS{S}{i}}, \m{\vpS{S}}, \m{\KcolS{S}},
\m{\BsysS{S}{\tau}}}{filling data and their components}
\nt{\m{\Trat{S}{R}}, \m{\Tlat{S}{R}}, \m{\GS{S}{R}}, \m{\GFin{S}}, \m{\latB{S}}, \m{\rfin{S}}, \m{\rinf{S}}, \m{\incH{S}},
\m{\fscale{S}{\tau}{a}}}{forms, groups, the lattice basis, the maps to the finite adeles, to infinity and from the
unitary group of the columns, frame scales}
\nt{\m{\KA{S}{u}}, \m{\KS{S}}, \m{\KgS{S}}, \m{\gid{S}}, \m{\redS{S}}, \m{\GRes{S}}, \m{\redim{S}{B}}, \m{\KT{S}{\sT}{u}}}{level
groups, the idele of $g$, reduction, the finite group, reduced image, level groups at~$\sT$}
\nt{$\psiS{S}$, \m{\charT{S}{\sT}}, \m{\uchar{S}{u}}, \m{\uord{S}}}{characters, unipotent character, the order}
\nt{\m{\Cl{S}}, $c$, $c'$, $y_c$, \m{\hc{S}{c}}, \m{\pos{S}{c}{c'}{\gamma}}, \m{\terms{S}{c}{s}{c'}}, \m{\clU{S}{c}}}{classes,
representatives, positions, terms, the unipotent lattice of a class}
\nt{\m{\Zf{S}{\gamma}}, \m{\tdisc{S}{\gamma}}, \m{\link{\gamma}}, $\Lambda$, \m{\tlink{S}{\Lambda}}}{frame Gram matrix,
term discriminant, links, the terms of a link}
\nt{\m{\Bc{S}}, $\sigma$, \m{\besfn{\sigma}{S}{y}{i}}, $W^{\sigma}$, \m{\bker{\sigma}{S}{\gamma}{y}}, \m{\cutoff{S}{c}},
\m{\amp{\sigma}{S}{c}{c'}{\gamma}}, \m{\ampabs{\sigma}{S}{c}{c'}{\gamma}}, \m{\ampsum{\sigma}{S}{\Xi}{c}{s}{c'}}}{Bessel
coordinates, a side, Bessel functions, kernel, cutoffs, amplitudes, amplitude sums}
\nt{\m{\intoff{S}{\sT}{x}}, \m{\loc{S}{\sT}}, \m{\isFL{S}{\sT}{\cL}}, \m{\cosamp{S}{\sT}{c}{s}{c'}{\gamma}}}{integral
off~$\sT$, components at~$\sT$, filling levels, coset amplitude}
\nt{$\cL$, \m{\gL{\cL}}, \m{\cgp{\cL}}, \m{\IsStd{\cL}}, \m{\gM{\cL}}, \m{\etab{j}}, \m{\fmod{\cL}{l}{g}},
\m{\umass{\cL}}, \m{\hdens{\cL}{\varrho}{p}{1}}}{a level, its unit, congruence subgroup, standard levels, group
measure, Heisenberg measures, floor modulus, unipotent mass, head density}
\nt{\m{\phs{g}{y}}, \m{\IntG{\cL}}, $Z$, \m{\Cos{Z}}, $a_c$, \m{\bvec{Z}{y}}, \m{\corr{Z}{p}{l}},
\m{\bcorr{\mu}{W}{\gamma}}}{finite Bessel function, intermediate groups, cosets, representatives, finite Bessel
vector, correlation, Bessel correlation}
\nt{$\gamma$, $\gamma_0$, $\gamma_1$, $\gamma_2$, \m{\gamma_{u,u'}}, \m{\gamma_{p,u,u'}}, $p$, $q$, $r$, $t$, $\omega$, $\zeta$, $d$}{rational
elements, sequences of units, the corner entry, the vector and coefficient of \m{\gamma_1z}, a product}
\nt{$k$, $l$, $l_0$, $l_1$, $x$, $\xi$, $P$, $u$, $u'$, $u_0$, $u_1$, $u_2$, $v'$, $w'$, $s$, $a$, $b$, $y$}{elements
of the level groups and adelic elements of the proofs, cosets, a point at infinity}
\nt{$L_p$, $\mathcal B$, $Q(u,u')$, \m{\mathsf r}, \m{\mathsf a_{u,u'}}, \m{\Sigma_c}, $C_s$, $k_s$, $B$, $\mu$, $J$,
$e$, $G_{\sT}$}{the limits, a block, the condition, integrands, sums of cutoffs, sets, representatives, a
subgroup, a measure, a Gram matrix, a difference, the group at~$\sT$}
\nt{$P$, $D$, \m{E(D)}, \m{m(D)}, \m{r(D)}, \m{k(D)}, \m{\varrho(D)}, \m{f(D)}, \m{b(D)}, \m{a(D)}, \m{N(D)}, \m{\hat\pi(D)},
\m{g_i(D)}, \m{\varpi(P)}, $N_0$}{parameters, admissible data and their numbers}
\nt{\m{X_0(D)}, \m{\Colv{D}{0}}, \m{\fs(D)}, \m{\sGD{D}}, \m{\SD{D}}, $K_f$, \m{\varkappa_D}, \m{\mathsf{w}^{\mathrm{col}}_D},
\m{\lambda}}{index set and level of the columns, the shape, the finite group, the filling data, integral points,
row parameter, column weight, the doubled parameter}
\nt{$\mathfrak B$, \m{\IsBes{\mathfrak B}{l}{w}}, \m{W^{\full}}, $A$, \m{\mathrm{Sch}}, \m{\rho}, \m{G_1(\fs)}}{a Bessel
system, its type, the full Bessel function, an intertwiner, the Schur representation, the real group}
\nt{\m{\kapS{S}{\tau}}, \m{c_{\mathrm{band}}}, \m{\bdeg{S}}, \m{\tailS{S}}, \m{\besE{S}}, \m{\mainE{S}}, \m{\codimF{S}},
\m{\packC{S}}, \m{\mass{E}{r}}}{constants of the set-up, the mass}
\nt{$\Delta$, $\Delta^{f,g}$, \m{x_{\tau,a,b}}, \m{\basisVec{D}{b}}, \m{\posEquiv{j}}, \m{\detRows{j}{i}{Y}{\tau}},
\m{\linePolyCoords{j}{i}{Y}}, $Y_i$, $L^{+}$, $P_l$, $p_l$, $f$, $g$, $n$, $j$, $i$, $i'$}{the objects of
Lemma~\ref{l:poly}}
\nt{$m$, $m'$, $\iota$, $n$, \m{\zeta_A}, \m{\zeta_{\hO}}, $\ell^2(X)$, $X$, $R$, $\Xi$, $N$, $K$, $F$, $O$}{numbers of a
level, zeta functions, \m{\CC^X}, a ring, a set of rational elements, and letters in names}
\end{multicols}
\endgroup

\end{document}
